 % 本文件由 scripts/assemble.py 自动装配，勿手改（改 parts/ 后重跑）
% 第2章 格波

% 第2章 LATTICE WAVES

\chapter{格波}

\epigraphbook{那些井然有序的运动，那规整齐一的步伐，虽不给耳朵任何声息，却向理解之心敲出最为盈满和谐的音律。}{Sir Thomas Browne}

\section{晶格动力学}

最简单的固体大概要数固态氩：它是中性原子的规则排列，原子带着结合紧密的封闭电子壳层，靠 van der Waals 力维系在一起，这种力主要作用在晶格中最近邻的原子之间。这类晶体的物理学，就是原子在各自理想化平衡位置附近的热运动。

要讨论这种运动，最初等的观念是 \emph{Einstein 模型}（Einstein model）：每个原子都在其邻居的力场所形成的势阱中，像一个简谐振子那样振动。这个力场绝不会严格是一个具有球对称性的二次势阱，但就平均而言，这是一种合理的近似。于是，晶体的激发谱由间距为 $\hbar\nu_E$ 的能级组成，其中 $\nu_E$ 是 \emph{Einstein 频率}（Einstein frequency），即每个原子在其势阱中的振荡频率。

这个模型在某些只需对热振动作极粗略处理的场合是有用的——尤其是在相对高温时，此时各原子彼此独立地振动这一假设是站得住脚的。但我们一眼就能看出：若两个或更多的原子一致地运动，它们之间那种趋于把每个原子拉回平衡位置的力便会减弱，因而激发一个量子所需的能量可能要小得多。相邻原子的运动有一种相互关联的趋势。

要得到整个晶格的谱，我们必须把局域的力放进来，并对运动作完整的描述。若不是晶格具有平移不变性，这将是一个无法解决的问题。这些激发必须满足 Bloch 定理，正如 §1.4 中所证明的那样。\emph{晶格动力学}（lattice dynamics）的理论，为第 1 章中建立起来的整套数学机器提供了最初等的例子。

让我们先定义\emph{晶格位移}（lattice displacements）：$\mathbf{u}_{sl}$ 是这样一个矢量，它给出第 $l$ 个晶胞中第 $s$ 个原子离开其平衡位置的位移。设此原子的质量为 $M_s$，则系统的动能定义为
\begin{equation*}
\text{K.E.} = \sum_{sl} \frac{1}{2} M_s \left| \dot{\mathbf{u}}_{sl} \right|^2.
\tag{2.1}
\end{equation*}
我们在这里作相当普遍的假设：晶格是带基元的格子（lattice with a basis），因此标记 $s$ 遍取晶胞中的一个或多个原子。

至于势能一项，我们本需要把原子间的力定义出来。但我们可以相当普遍地假设，存在一个函数 $\mathcal{V}(\mathbf{u}_{sl})$，它用所有原子的位置——也就是说，用它们在某一时刻离开平衡晶格格点的实际位移——来表达整个晶体的势能。我们还可以进一步假设：当所有 $\mathbf{u}_{sl}$ 为零时，这个函数取极小，因为完整晶格想必是一个稳定平衡的组态。我们利用微振动理论，把 $\mathcal{V}$ 在这一点附近按 $\mathbf{u}_{sl}$ 的笛卡儿分量 $u_{sl}^{j}$ 展开成幂级数：
\begin{equation*}
\mathcal{V} = \mathcal{V}_0 + \sum_{slj} u_{sl}^{j} \left[ \frac{\partial \mathcal{V}}{\partial u_{sl}^{j}} \right]_0
+ \frac{1}{2} \sum_{ss',\,ll',\,jj'} u_{sl}^{j} u_{s'l'}^{j'} \left[ \frac{\partial^2 \mathcal{V}}{\partial u_{sl}^{j} \, \partial u_{s'l'}^{j'}} \right]_0 + \ldots
\tag{2.2}
\end{equation*}
常数项在当前的讨论里无关紧要；线性项的系数必须为零，因为我们处在平衡点附近；于是第一个重要的项必然是位移的二次项——即所谓的\emph{谐和项}（harmonic term）。

由这两个函数，即式 (2.1) 与式 (2.2)，用经典力学中通常的 Lagrange 程序，我们得到各分量的运动方程
\begin{equation*}
M_s \ddot{u}_{sl}^{j} = - \sum_{s'l'} \left[ \frac{\partial^2 \mathcal{V}}{\partial u_{sl}^{j} \, \partial u_{s'l'}^{j'}} \right]_0 u_{s'l'}^{j'}
\tag{2.3}
\end{equation*}
它对晶胞中所有的格点 $s$、对晶格中所有的晶胞 $l$、以及位移矢量的每一个笛卡儿分量 $j = x, y, z$ 都成立。

这些方程当然代表了数目极其庞大的耦合微分方程。但让我们考察一下右边的系数。令
\begin{equation*}
\left[ \frac{\partial^2 \mathcal{V}}{\partial u_{sl}^{j} \, \partial u_{s'l'}^{j'}} \right]_0 \equiv \mathsf{G}_{sl;s'l'}^{jj'}.
\tag{2.4}
\end{equation*}
这些系数构成一个笛卡儿张量的分量，于是我们可以把式 (2.3) 用矢量形式写得更简洁些：
\begin{equation*}
M_s \ddot{\mathbf{u}}_{sl} = - \sum_{s'l'} \mathsf{G}_{sl;s'l'} \cdot \mathbf{u}_{s'l'}.
\tag{2.5}
\end{equation*}

这个方程可以作如下解释：右边求和中的每一项，都是第 $l'$ 个晶胞中第 $s'$ 个格点上的原子因位移 $\mathbf{u}_{s'l'}$ 而\emph{作用在第 $l$ 个晶胞中第 $s$ 个原子上的力}。倘若我们知道整个晶格的势能可以由原子对之间的力构造出来，那就可以直接用原子间力常数把它写出来。

但无论这些力的物理图像如何，有一点是必需的：它们不能依赖于 $l$ 和 $l'$ 在晶格中的绝对位置。张量 $\mathsf{G}$ 必须只是它们相对位置的函数：
\begin{equation*}
\mathsf{G}_{sl;s'l'} = \mathsf{G}_{ss'}(\boldsymbol{h}), \quad \text{其中}\ \ \boldsymbol{h} = \boldsymbol{l}' - \boldsymbol{l}.
\tag{2.6}
\end{equation*}
我们的运动方程现在读作
\begin{equation*}
M_s \ddot{\mathbf{u}}_{sl} = - \sum_{s'\boldsymbol{h}} \mathsf{G}_{ss'}(\boldsymbol{h}) \cdot \mathbf{u}_{s,\,l+\boldsymbol{h}}.
\tag{2.7}
\end{equation*}

这些方程具有满足 Bloch 定理所要求的平移不变的形式；把标记从 $l$ 换成譬如 $l''$，只不过是原封不动地重复同一组方程。现在假设我们已经找到了一个解——一组描述各个 $l$ 值的 $\mathbf{u}_{sl}$ 的时间函数。这些函数必须满足 Bloch 条件 (1.48) 或 (1.57)。于是存在一个波矢\footnote{我们用这个符号表示格波的波矢，而用 $\boldsymbol{k}$ 表示电子波的波矢。当然，二者都是同一倒格子空间中的矢量。}$\boldsymbol{q}$，使得
\begin{equation*}
\mathbf{u}_{sl}(t) = e^{i\boldsymbol{q}\cdot\boldsymbol{l}} \mathbf{u}_{s,o}(t),
\tag{2.8}
\end{equation*}
其中 $\mathbf{u}_{s,o}(t)$ 是碰巧被选作格矢 $\boldsymbol{l}$ 原点的那一晶胞处的位移。注意：在每一个晶胞里，格点 $s$ 上的原子都以相同的振幅沿相同的方向运动；逐胞不同的只是相位。

动力学方程的每一个解，都有一个使式 (2.8) 得以满足的 $\boldsymbol{q}$ 值。要弄清哪个解带哪个波矢，我们必须回到式 (2.7)。把式 (2.8) 代入，得
\begin{equation*}
M_s \ddot{\mathbf{u}}_{s,o}\, e^{i\boldsymbol{q}\cdot\boldsymbol{l}} = - \sum_{s'\boldsymbol{h}} \mathsf{G}_{ss'}(\boldsymbol{h}) \cdot \mathbf{u}_{s',o}\, e^{i\boldsymbol{q}\cdot\boldsymbol{h}} e^{i\boldsymbol{q}\cdot\boldsymbol{l}}.
\tag{2.9}
\end{equation*}
因子 $\exp i\boldsymbol{q}\cdot\boldsymbol{l}$ 可以消去。为了表明原点的选取完全任意，而我们处理的乃是具有确定 $\boldsymbol{q}$ 值的一个解，让我们写
\begin{equation*}
\mathbf{u}_{s,o} = \mathbf{U}_{s,\boldsymbol{q}}.
\tag{2.10}
\end{equation*}

于是只剩下
\begin{align*}
M_s \ddot{\mathbf{U}}_{s,\boldsymbol{q}} &= - \sum_{s'} \Bigl\{ \sum_{\boldsymbol{h}} \mathsf{G}_{ss'}(\boldsymbol{h})\, e^{i\boldsymbol{q}\cdot\boldsymbol{h}} \Bigr\} \cdot \mathbf{U}_{s',\boldsymbol{q}} \\
&= - \sum_{s'} \mathsf{G}_{ss'}(\boldsymbol{q}) \cdot \mathbf{U}_{s',\boldsymbol{q}},
\tag{2.11}
\end{align*}
其中
\begin{equation*}
\mathsf{G}_{ss'}(\boldsymbol{q}) \equiv \sum_{\boldsymbol{h}} \mathsf{G}_{ss'}(\boldsymbol{h})\, e^{i\boldsymbol{q}\cdot\boldsymbol{h}}
\tag{2.12}
\end{equation*}
——它是力张量（force tensor）$\mathsf{G}$ 的一个 Fourier 变换。

新方程组 (2.11) 相对而言容易求解。关键在于它的数目是有限的。设每个晶胞含 $n$ 个原子，晶格共有 $N$ 个晶胞。原来的方程组 (2.3) 或 (2.5) 将有 $3nN$ 个——这里对矢量的三个笛卡儿分量是分别计算的。而新方程组 (2.11) 只有 $3n$ 个——$n$ 个 $s$ 值的每一个各有 3 个分量方程。

这是平移不变性的一个极有威力的例证。所有晶胞在结构上本质同一，这意味着我们只需要关于其中一个晶胞的信息，就能计算整个集合的动力学。诚然，若想为每个 $\boldsymbol{q}$ 值（边界条件恰好允许 $N$ 个不同的值）求得解，我们得分别计算；但在实践中我们只关心这类解的一个有限样本，其余的可以靠插值得到。

接下来的步骤是初等的。如同在通常的振动理论中那样，我们假设 $\mathbf{U}_{s,\boldsymbol{q}}$ 含有时间因子 $\exp(i\nu t)$。于是必须求解由 $3n$ 个方程组成的方程组
\begin{equation*}
\sum_{s'j'} \Bigl\{ \mathsf{G}_{ss'}^{jj'}(\boldsymbol{q}) - \nu^2 M_s \delta_{ss'} \delta_{jj'} \Bigr\} U_{s',\boldsymbol{q}}^{j'} = 0,
\tag{2.13}
\end{equation*}
其未知对象是 $\mathbf{U}_{s,\boldsymbol{q}}$ 的各分量。这是一个本征值方程；令矩阵 $\{\ \}$ 的行列式为零，得到一个关于 $\nu^2$ 的方程，求出它的根，便得到全部 $3n$ 个解。实际上，我们所做的无非是求出晶胞内 $n$ 个原子的 $3n$ 个简正模（normal mode）振动，只是假定它们通过力张量 $\mathsf{G}_{ss'}(\boldsymbol{q})$ 相互作用。这个张量不同于原来的原子间力张量 $\mathsf{G}_{ss'}(\boldsymbol{h})$，而且对每个 $\boldsymbol{q}$ 值都不相同；对给定的 $s$ 和 $s'$，它把晶格中每个晶胞的格点 $s$ 上的\emph{所有} $s$ 类原子（即每个晶胞中位于格点 $s$ 上的那些原子）与格点 $s'$ 上的\emph{所有}原子之间的相互作用总括起来，并计及它们之间的相对相位。

\section{格波的性质}

要理解格波的物理性质，研究几个简单的例子是有益的。其中最简单的当属\emph{线性链}（linear chain）。

假设我们有一维“晶格”，每个晶胞一个原子，且只有最近邻力。运动方程很容易写出来。势能（参见式 (2.2)）具有如下形式：
\begin{equation*}
\mathcal{V} = \sum_{l} \frac{1}{2} \alpha \left( u_{l} - u_{l+a} \right)^2,
\tag{2.14}
\end{equation*}
其中 $\alpha$ 是力常数（force constant）。运动方程（参见式 (2.5)）变为
\begin{equation*}
M \ddot{u}_{l} = - \alpha \left( 2u_{l} - u_{l+a} - u_{l-a} \right)
\tag{2.15}
\end{equation*}

\begin{figure}[H]
\centering
\includegraphics[width=0.72\textwidth]{fig_16.pdf}
\caption*{图 16\quad 线性链的振动频率}
\end{figure}

通过代换
\begin{equation*}
u_{l} = U_{q}\, e^{iql},
\tag{2.16}
\end{equation*}
它变换成与式 (2.11) 类比的方程，即
\begin{align*}
M \ddot{U}_{q} &= - \left( 2\alpha - \alpha e^{iqa} - \alpha e^{-iqa} \right) U_{q} \\
&= - 2\alpha \left( 1 - \cos qa \right) U_{q},
\tag{2.17}
\end{align*}
其中 $a$ 是链中原子的间距。这是一个简谐振子的方程，其频率为
\begin{equation*}
\nu_{q} = \sqrt{\left( \frac{\alpha}{M} \right)}\, 2 \sin \left( \frac{qa}{2} \right).
\tag{2.18}
\end{equation*}

这个众所周知的结果体现了理论的许多要点。（i）一切可能的振动，由如下范围内的 $q$ 值给出：
\begin{equation*}
- \frac{\pi}{a} < q \leqslant \frac{\pi}{a}.
\tag{2.19}
\end{equation*}
这正是我们的布里渊区（Brillouin zone）。任何落在这个范围之外的 $q$ 值都只是精确地重复同一运动：
\begin{equation*}
u_{l} = U\, e^{i(q+\sigma)l} \equiv U\, e^{iql}.
\tag{2.20}
\end{equation*}
把波矢取在布里渊区之外是相当不自然的（尽管并非没有意义！）。不过，假如我们当初用了大于 $\pi/a$ 的 $|q|$ 值，其实也会得到正确的结果。在式 (2.18) 中，$\nu_q$ 是 $q$ 的周期函数，所以凡在布里渊区内约化到同一点的那些 $q$ 值都具有相同的频率。

恰好有 $N$ 个不同的解，对应于布里渊区中 $N$ 个允许的 $q$ 值。这与原晶格的 $N$ 个自由度正相一致。

（ii）对于小的 $q$ 值（即 $qa \ll 1$），
\begin{equation*}
\nu_{q} \sim \sqrt{\left( \frac{\alpha}{M} \right)}\, aq.
\tag{2.21}
\end{equation*}

频率与波数成正比，这对应于连续介质中普通弹性波的一个众所周知的性质：若扰动的波长远大于晶格常数，那么这条原子链在经典力学中的行为就像一根“重弹性绳”。

然而在大的 $q$ 值处，波的速度不再是常数；到了 $q = \pi/a$，即波长等于 $2a$ 时，函数 $\nu_q$ 弯折下来趋于一条水平切线。这表现出\emph{色散}（dispersion）性质。

现在考虑一个更复杂的情形：一条原子线性链，间距和力常数都与前面相同，但两种不同的质量 $M_1$ 与 $M_2$ 交替排列。现在我们的晶胞含两个原子。与式 (2.11) 类比的方程要比式 (2.17) 复杂：
\begin{equation*}
\left.
\begin{aligned}
M_1 \ddot{U}_{1} &= - 2\alpha U_{1} + 2\alpha \cos qa \cdot U_{2}, \\
M_2 \ddot{U}_{2} &= - 2\alpha U_{2} + 2\alpha \cos qa \cdot U_{1}.
\end{aligned}
\right\}
\tag{2.22}
\end{equation*}

为求频率 $\nu$，我们必须求解行列式方程
\begin{equation*}
\begin{vmatrix}
2\alpha - M_1 \nu^2 &\quad - 2\alpha \cos qa \\
- 2\alpha \cos qa &\quad 2\alpha - M_2 \nu^2
\end{vmatrix} = 0,
\tag{2.23}
\end{equation*}
它有两个根
\begin{equation*}
\nu_{\pm}^{2} = \alpha \left( \frac{1}{M_1} + \frac{1}{M_2} \right) \pm \alpha \sqrt{\left\{ \left( \frac{1}{M_1} + \frac{1}{M_2} \right)^{2} - \frac{4 \sin^{2} qa}{M_1 M_2} \right\}}.
\tag{2.24}
\end{equation*}
把这两个根作为 $q$ 的函数描画出来，就如图 18 所示。

\begin{figure}[H]
\centering
\includegraphics[width=0.72\textwidth]{fig_17.pdf}
\caption*{图 17\quad 双原子线性链}
\end{figure}

\begin{figure}[H]
\centering
\includegraphics[width=0.72\textwidth]{fig_18.pdf}
\caption*{图 18\quad 双原子链的振动频率}
\end{figure}

与单原子链一样，这里有一个根 $\nu_{-}$，它在 $q = 0$ 附近趋于与 $q$ 成正比。我们称之为\emph{声学支}（acoustic mode），因为它是把链当作弹性连续体时的长波长振动的对应物。这种情形下求出的声速正是这种介质中正常的声速——我们把这个证明留作读者的练习。

但还有另一个分支 $\nu_{+}$，在 $q = 0$ 附近满足
\begin{equation*}
\nu_{+}^{2} \sim 2\alpha \left( \frac{1}{M_1} + \frac{1}{M_2} \right)
\tag{2.25}
\end{equation*}
这个分支与声学支相隔甚远，但随着 $q$ 的增大，两个分支的频率有相互靠近的趋势。这叫做\emph{光学支}（optical mode）。我们容易验证：在 $q = 0$ 的情形，轻、重两类原子各自的子晶格（sublattice）都在作刚体式的整体运动，且两者彼此反向；或者可以这样说，每个晶胞里的双原子分子都在振动，仿佛与邻居无关。如果像在离子晶体中那样，两类原子带有异号电荷，这就给出一个振荡的偶极矩，因而是光学活性的。

\begin{figure}[H]
\centering
\includegraphics[width=0.72\textwidth]{fig_19.pdf}
\caption*{图 19\quad 振动频率：(a) 按双原子链处理的单原子线性链；(b) 单原子线性链，显示区边界；(c) 按单原子链处理的双原子线性链}
\end{figure}

研究从图 16 的单原子链到图 18 的双原子链的转变，是很有启发性的。假设我们在第一种情形下一开始就取一个含两个原子的晶胞，而无视这两个原子质量相同的事实。容易证明：此时光学支与声学支将在 $q = \pi/2a$ 处相接，如图 19(a) 所示。与图 16 相比，我们似乎把模的数目加倍了。然而，晶胞长度加倍之后，布里渊区的尺寸显然减半了。通过平移 $\pi/a$（即对我们这个加倍后的晶胞而言的一个倒格矢），可以把图 19(a) 展开成图 19(b)。这时“声学”支与“光学”支便连续相接。图 19(a) 的简约区是虚假地只及应有大小的一半。

我们可以换一种说法。交替改变链上原子的质量，其效果是在 $\pm \pi/2a$ 处引入新的区边界。频率跨越这条边界不再连续；出现了一个隙。我们可以用图 19(c) 来表示——这是处理双原子链的另一种方式。在下一章里我们会经常遇到这个现象。

下一个要考虑的情形，该是具有最近邻相互作用的二维或三维 Bravais 格子。然而这时方程已经复杂得难以轻松地一般求解。每一个维度都给位移矢量引入一个额外的笛卡儿分量，因此在三维情形我们必须求解一个关于 $\nu^2$ 的三次方程。

在这种情形下，只要过渡到长波极限，这三个根就可以在物理上加以辨认。实际上我们处理的是一个弹性连续体，其结构过于细微，在长波的动力学中无足轻重。众所周知，在这样的连续体中可以传播三种不同类型、速度各不相同的声波。

\begin{figure}[H]
\centering
\includegraphics[width=0.72\textwidth]{fig_20.pdf}
\caption*{图 20\quad 恒定频率面：(a) 各向同性介质；(b) 真实晶体}
\end{figure}

这三个\emph{声学支}的本质差别在于它们的\emph{偏振}（polarization）性质。例如在各向同性连续体中，有一个模是\emph{纵}偏振的——每个原子的位移矢量都沿着波的传播方向。另外两个模速度相同，是\emph{横}偏振的——原子在垂直于波矢的平面内运动。一般说来，纵模的速度比横模高，因此 $\boldsymbol{q}$ 空间中的恒定频率面看上去就像图 20(a) 那样。

晶体固体在宏观弹性性质上通常不是各向同性的，所以这幅图像会引人误解。给定类型的波，其速度将依赖于传播方向。于是（图 20(b)）横支的恒定频率面可能离球面很远。这样，除了特殊的对称方向之外，两个横模将不再简并，它们的频率面可能以复杂的样式相互交截。的确，把模形式上分成横模与纵模是有点误导的：在这样的介质中，实际的偏振矢量（即方程 (2.11) 或 (2.13) 的解矢量 $\mathbf{U}_{\boldsymbol{q}}$ 的方向）未必严格地平行或垂直于 $\boldsymbol{q}$。

由此，格波的长波行为就确定下来了。对沿给定方向的 $\boldsymbol{q}$，我们必须有 $\nu$ 与 $q$ 之间的一个关系，其形式如图 21 所勾勒——在 $q = 0$ 附近是线性关系，斜率等于相应的宏观声速。但随着 $q$ 增大，我们逼近布里渊区的边界，就必定出现色散。

\begin{figure}[H]
\centering
\includegraphics[width=0.72\textwidth]{fig_21.pdf}
\caption*{图 21\quad 格波在给定方向上的色散}
\end{figure}

此后的确切行为会相当复杂。一般地说，$\nu_{\boldsymbol{q}}$ 会像线性情形那样弯折下来，但不一定弯到水平切线。唯一的规则是：$\nu_{\boldsymbol{q}}$ 作为倒格子空间中 $\boldsymbol{q}$ 的函数，其根的整体图样是连续的，并且具有倒格子的周期性。察看式 (2.8) 与式 (2.12) 即可明白这一点。与线性情形一样，若 $\boldsymbol{g}$ 是一个倒格矢，则波矢 $\boldsymbol{q}' = \boldsymbol{q} + \boldsymbol{g}$ 的解与波矢 $\boldsymbol{q}$ 的解完全相同，所以 $\nu_{\boldsymbol{q}}$ 以 $\boldsymbol{g}$ 为周期。它是 $\boldsymbol{q}$ 空间中的连续函数，因为它是一个本征值方程的解，该方程的矩阵是对称的，而其系数本身就是 $\boldsymbol{q}$ 空间中的连续周期函数。

这意味着，当我们把像图 20(b) 那样的图形扩展到整个区时，它必定构成这种图样的一部分，如图 22 所草示。只要知道力常数，这样的图是能够计算出来的——尽管在实践中更常见的做法是：先用中子或 X 射线衍射（参见 §2.8）测定各方向上的色散曲线，再反过来寻找能与之拟合的原子力常数。

在带基元的三维晶格这一更复杂的情形，我们还会发现横光学支与纵光学支，它们各有自己的色散曲线和能量面。所有这些情形，原则上都涵盖在我们的普遍方程 (2.13) 之中。

%（以下为 §2.2 收尾的插图，印于书页 37 顶部，随图移归上一部分 A）
\begin{figure}[H]
\centering
\includegraphics[width=0.72\textwidth]{fig_22.pdf}
\caption*{图 22\quad 重复区图式（repeated zone scheme）下晶格频率函数的一支 $\nu_{\mathbf{q}}$}
\end{figure}

\section{晶格求和}

在长波极限下，晶格简正模方程的解与宏观连续介质弹性波方程的解之间的等价性，为把原子力常数同固体中观测到的弹性常数联系起来提供了一个便利的方法。原则上这是一个比较简单的纲领，尽管在实践中它可能涉及繁冗的计算。

有一个主要的困难之处。在一个常见的情形中，我们对原子间力了解得非常清楚：在离子晶体中，固体内聚能的主要部分来自异号电荷离子之间的 Coulomb 力。例如，对于位于诸点 $\mathbf{R}_l$ 处的一组离子，总静电能必定由
\begin{equation*}
\mathcal{E} = \frac{1}{2}\sum_{ij} \frac{\pm e^2}{|\mathbf{R}_i - \mathbf{R}_j|}
\tag{2.26}
\end{equation*}
给出，正负号取决于 $\mathbf{R}_i$ 与 $\mathbf{R}_j$ 两处电荷的相对符号。

此式可以改写成
\begin{equation*}
\mathcal{E} = -\frac{1}{2}N\alpha\,\frac{e^2}{a},
\tag{2.27}
\end{equation*}
其中 $a$ 是一个晶格常数，系数 $\alpha$ 现在是无量纲的，依赖于正、负离子在晶格中的排布方式。这样，\emph{Madelung 常数} $\alpha$ 作为表征某一特定类型晶体结构的参量（而不管其尺度大小），就有了一定的意义。

\begin{figure}[H]
\centering
\includegraphics[width=0.72\textwidth]{fig_23.pdf}
\caption*{图 23\quad （原书此图无图注）}
\end{figure}

例如，设我们的离子晶体由两套子晶格组成——正离子位于一组格点 $\boldsymbol{l}$ 上，负离子位于 $\boldsymbol{l}+\mathbf{x}$ 上。我们需要计算
\begin{equation*}
\alpha = a\sum_{\boldsymbol{l}}\left\{\frac{1}{|\boldsymbol{l}|} - \frac{1}{|\boldsymbol{l}+\mathbf{x}|}\right\}
\tag{2.28}
\end{equation*}
（第一项中除去 $\boldsymbol{l}=0$）。麻烦在于，这个和式中的两项分开来各自发散；在距离 $\boldsymbol{l}$ 处，有待计入的点的数目是 $l^2$ 的量级。

这个级数是条件收敛的，而且即使如此，收敛也很慢。计算这类和式的问题变得十分严峻。最好的直接方法是 Evjen 方法，即从原点向外逐个处理各层壳，使每一壳的总静电电荷恰好为零。

这个困难并不只限于内聚能的计算。在计算诸如式 (2.12) 那样的表达式——原子之间力张量的 Fourier 变换——时也会出现。在那里，我们同样要对所有格点求和，而对 Coulomb 力来说收敛很慢。任何计算离子晶体晶格动力学的尝试都必然碰上这个困难。

有一个非常漂亮的方法，称为 \emph{Ewald 方法}（Ewald's method），可以用它解决这个问题。也许这只是一个技巧，算不上固体理论的一条基本原理，但它实在精巧，值得稍作题外的讨论。它还演示了晶格上周期函数理论的某些内容。

先考虑函数
\begin{equation*}
\frac{2}{\sqrt{\pi}}\sum_{\boldsymbol{l}} e^{-|\boldsymbol{l}-\mathbf{r}|^2\rho^2} = F(\mathbf{r},\rho).
\tag{2.29}
\end{equation*}
这个函数在 $\mathbf{r}$ 中是周期的，具有晶格的周期性。因此，按式 (1.15)，可以把它展开成 Fourier 级数
\begin{equation*}
F(\mathbf{r},\rho) = \sum_{\boldsymbol{g}} F_{\boldsymbol{g}}\, e^{i\boldsymbol{g}\cdot\mathbf{r}},
\tag{2.30}
\end{equation*}
其中
\begin{equation*}
F_{\boldsymbol{g}} = \frac{1}{V}\int \frac{2}{\sqrt{\pi}}\sum_{\boldsymbol{l}} e^{-|\boldsymbol{l}-\mathbf{r}|^2\rho^2}\, e^{-i\boldsymbol{g}\cdot\mathbf{r}}\, d\mathbf{r},
\tag{2.31}
\end{equation*}
如式 (1.22)（积分遍及整个晶体）。

现在我们在级数的每一项中引入一个因子 $\exp(i\boldsymbol{g}\cdot\boldsymbol{l})=1$。这样一来，
\begin{align*}
F_{\boldsymbol{g}} &= \frac{1}{V}\,\frac{2}{\sqrt{\pi}}\sum_{\boldsymbol{l}}\int e^{-|\boldsymbol{l}-\mathbf{r}|^2\rho^2}\, e^{-i\boldsymbol{g}\cdot(\mathbf{r}-\boldsymbol{l})}\, d\mathbf{r} \\
&= \frac{N}{V}\,\frac{2}{\sqrt{\pi}}\int e^{-r^2\rho^2-i\boldsymbol{g}\cdot\mathbf{r}}\, d\mathbf{r},
\tag{2.32}
\end{align*}
因为和式中的每一项现在都相同，只不过各自取晶格中不同的点 $\boldsymbol{l}$ 为原点来计算。由于积分收敛很快，而所有格点又彼此等价，这一步是显而易见的。

式 (2.32) 中的积分很容易算出。它给出
\begin{equation*}
F_{\boldsymbol{g}} = \frac{2\pi}{v_c}\,\frac{1}{\rho^3}\, e^{-g^2/4\rho^2}.
\tag{2.33}
\end{equation*}
把式 (2.33) 代入式 (2.29) 和 (2.30)，得到的关系式在纯数学中称为 \emph{$\theta$ 函数变换}（theta-function transformation）（在纯数学中通常只对简单立方格子证明这个关系；我们的结果更为普遍）
\begin{equation*}
\frac{2}{\sqrt{\pi}}\sum_{\boldsymbol{l}} e^{-|\boldsymbol{l}-\mathbf{r}|^2\rho^2} = \frac{2\pi}{v_c}\sum_{\boldsymbol{g}}\frac{1}{\rho^3}\, e^{-g^2/4\rho^2}\, e^{i\boldsymbol{g}\cdot\mathbf{r}}.
\tag{2.34}
\end{equation*}
但是考虑恒等式
\begin{equation*}
\frac{2}{\sqrt{\pi}}\int_0^\infty e^{-z^2\rho^2}\, d\rho = \frac{1}{|z|}.
\tag{2.35}
\end{equation*}
把它用到式 (2.34) 的左端，我们有
\begin{equation*}
\sum_{\boldsymbol{l}}\frac{1}{|\boldsymbol{l}-\mathbf{r}|} = \sum_{\boldsymbol{l}}\frac{2}{\sqrt{\pi}}\int_0^\infty e^{-|\boldsymbol{l}-\mathbf{r}|^2\rho^2}\, d\rho.
\tag{2.36}
\end{equation*}
我们把对哑变量 $\rho$ 的积分在某个任意点 $G$ 处分成两部分，并在第一部分中用式 (2.34) 代换被积函数：
\begin{align*}
\sum_{\boldsymbol{l}}\frac{1}{|\boldsymbol{l}-\mathbf{r}|} &= \frac{2\pi}{v_c}\int_0^G \sum_{\boldsymbol{g}}\frac{1}{\rho^3}\, e^{-g^2/4\rho^2}\, e^{i\boldsymbol{g}\cdot\mathbf{r}}\, d\rho + \sum_{\boldsymbol{l}}\frac{2}{\sqrt{\pi}}\int_G^\infty e^{-|\boldsymbol{l}-\mathbf{r}|^2\rho^2}\, d\rho \\
&= \frac{\pi}{v_c}\,\frac{1}{G^2}\sum_{\boldsymbol{g}} e^{i\boldsymbol{g}\cdot\mathbf{r}}\,\frac{e^{-g^2/4G^2}}{g^2/4G^2} + \sum_{\boldsymbol{l}}\frac{1}{|\boldsymbol{l}-\mathbf{r}|}\,\mathrm{erfc}\{G|\boldsymbol{l}-\mathbf{r}|\},
\tag{2.37}
\end{align*}
其中第二项含有\emph{余误差函数}（complementary error function），当自变量取大值时它迅速趋于零。

具体做法是选取一个 $G$ 值，使得对 $\boldsymbol{l}$ 的级数和对 $\boldsymbol{g}$ 的级数都迅速收敛；容易看出，这两个级数都比原来的级数 (2.36) 收敛得好得多。从“物理上”讲，这就好像我们分两步在 $\mathbf{r}$ 处建立起电荷。第一步，我们作一个球对称的 Gaussian 分布，其在半径 $z$ 处的电荷密度正比于 $\exp(-G^2z^2)$。点电荷晶格与这个电荷云在长距离上的静电相互作用，给出式 (2.37) 中的第一项。然后我们再作修正：在 $\mathbf{r}$ 处放一个点电荷，同时减去 Gaussian 分布。这个对象是电中性的，它与晶格的长程相互作用可以忽略，只贡献一个对邻近格点迅速收敛的和。

不过事情还没有完，因为上面的和式全部取正号，而且当 $\mathbf{r}=0$ 时它包含了原点处的电荷对自身的作用。这会表现为第二个求和中 $\boldsymbol{l}=0$ 那一项在 $\mathbf{r}=0$ 附近的一个发散贡献。我们从左端的级数中减去 $1/r$，这等价于把右端对 $\boldsymbol{l}$ 求和中的第一项去掉，而在 $r\to 0$ 的极限下代之以下式
\begin{align*}
\frac{1}{r}\,\mathrm{erf}\,Gr - \frac{1}{r} &= \frac{2}{\sqrt{\pi}}\,\frac{1}{r}\left\{\int_{Gr}^{\infty} e^{-z^2}\, dz - \frac{\sqrt{\pi}}{2}\right\} \\
&= -\frac{2}{\sqrt{\pi}}\,\frac{1}{r}\int_0^{Gr} e^{-z^2}\, dz \\
&\to -2G/\sqrt{\pi}.
\tag{2.38}
\end{align*}
于是
\begin{equation*}
\sum_{\boldsymbol{l}\neq 0}\frac{1}{|\boldsymbol{l}|} = \frac{\pi}{v_c}\,\frac{1}{G^2}\sum_{\boldsymbol{g}}\frac{e^{-g^2/4G^2}}{g^2/4G^2} + \sum_{\boldsymbol{l}\neq 0}\frac{1}{|\boldsymbol{l}|}\,\mathrm{erf}\{G|\boldsymbol{l}|\} - \frac{2G}{\sqrt{\pi}}.
\tag{2.39}
\end{equation*}
式 (2.37) 和 (2.39) 都还含有一个奇异性；$\boldsymbol{g}=0$ 的项是发散的。但当我们像式 (2.28) 那样计算 Madelung 常数时，是从一个级数中减去另一个级数，两个奇异性便相互抵消。这说白了就是：每一套子晶格各自的平均势如果单独计算将是无穷大，但晶体的总电荷为零。

从这里我们可以看到这一变换的某种物理意义。考虑一个更复杂的情形。设有一个格波，在晶格的每个格点上产生一个偶极子，
\begin{equation*}
\mathbf{p}(\boldsymbol{l}) = \mathbf{p}\, e^{i\mathbf{q}\cdot\boldsymbol{l}}.
\tag{2.40}
\end{equation*}
由初等静电学可知，这组电荷分布在 $\mathbf{r}$ 处产生的 Coulomb 场，就是这些偶极子所产生的势的梯度：
\begin{align*}
\mathbf{E}(\mathbf{r}) &= \nabla\left\{\sum_{\boldsymbol{l}} \mathbf{p}(\boldsymbol{l})\cdot\nabla\left(\frac{1}{|\boldsymbol{l}-\mathbf{r}|}\right)\right\} \\
&= \nabla(\mathbf{p}\cdot\nabla)\left\{\sum_{\boldsymbol{l}} \frac{e^{i\mathbf{q}\cdot\boldsymbol{l}}}{|\boldsymbol{l}-\mathbf{r}|}\right\}.
\tag{2.41}
\end{align*}
对格点的这个和式是式 (2.36) 稍微复杂一点的版本。用我们上面在 $\mathbf{q}=0$ 的特殊情形中用过的同一套论证就可以把它求出来；结果是
\begin{equation*}
\sum_{\boldsymbol{l}} \frac{e^{i\mathbf{q}\cdot\boldsymbol{l}}}{|\boldsymbol{l}-\mathbf{r}|} = \frac{\pi}{v_c}\,\frac{1}{G^2}\sum_{\boldsymbol{g}} \frac{e^{i(\mathbf{q}+\boldsymbol{g})\cdot\mathbf{r}}\, e^{-|\mathbf{q}+\boldsymbol{g}|^2/4G^2}}{|\mathbf{q}+\boldsymbol{g}|^2/4G^2} + \sum_{\boldsymbol{l}} \frac{e^{i\mathbf{q}\cdot\boldsymbol{l}}}{|\boldsymbol{l}-\mathbf{r}|}\,\mathrm{erf}\{G|\boldsymbol{l}-\mathbf{r}|\},
\tag{2.42}
\end{equation*}
它可以逐项微商，正是式 (2.41) 所需要的。事实上，当离子之间存在 Coulomb 力时，我们在这里得到的就是计算力张量系数（如式 (2.12)）所需的那类公式。

但是考虑第一个求和中 $\boldsymbol{g}=0$ 那一项对电场的贡献：
\begin{equation*}
\mathbf{E}_0(\mathbf{r}) = \frac{-4\pi}{v_c}\, e^{-q^2/4G^2}\,\frac{(\mathbf{p}\cdot\mathbf{q})}{q^2}\,\mathbf{q}.
\tag{2.43}
\end{equation*}
除了因子 $\exp(-q^2/4G^2)$ 之外，这正是宏观电场 $\overline{\mathbf{E}}(\mathbf{r})$，即与下述宏观极化波
\begin{equation*}
\mathbf{p}(\mathbf{r}) = \mathbf{p}\, e^{i\mathbf{q}\cdot\mathbf{r}}
\tag{2.44}
\end{equation*}
相联系的场——此时把晶体当作连续介质看待。于是我们可以把式 (2.41) 表成
\begin{equation*}
\mathbf{E}(\mathbf{r}) = \overline{\mathbf{E}}(\mathbf{r}) + \frac{4\pi}{v_c}\left(1-e^{-q^2/4G^2}\right)\frac{(\mathbf{p}\cdot\mathbf{q})}{q^2}\,\mathbf{q} + \sum_{\boldsymbol{g}\neq 0}\{\,\} + \sum_{\boldsymbol{l}}\{\,\},
\tag{2.45}
\end{equation*}
其中我们不必费神写出对 $\boldsymbol{g}$ 和对 $\boldsymbol{l}$ 求和的所有各项，只须注意它们全都迅速收敛（$\mathbf{r}=0$ 的邻域除外；在该邻域可以采用式 (2.39) 中用过的办法消除原点处偶极子的自能）。在这个表达式中，我们把在 $\mathbf{q}=0$ 附近表现很坏的宏观场 $\overline{\mathbf{E}}(\mathbf{r})$ 分离了出来，其余各项则对一切 $\mathbf{q}$ 值都迅速收敛于一个确定的局部场（local field）（参见 §8.2）。

这个局部场在离子晶体的晶格动力学中是重要的，因为它可能使离子本身发生极化。由此产生的感生偶极子分布又反过来对作用于离子上的力作出贡献，因此必须包括在有效力张量（式 (2.4)）之中：点电荷之间的中心成对力不足以解释实验事实。从这个方向入手，理论公式显得相当繁复，但它们可以从一个简单的物理模型导出：把每个离子换成一个质量较大、电荷设为 $e^+$ 的“芯”，它连着一个无重量、由电荷为 $e^-$ 的价电子组成的“壳层”。这个\emph{壳层模型}（shell model）中的各种电荷和弹性常数，不仅提供了足以把晶格谱（lattice spectrum）拟合到实验上的可调参数，而且解析地表示了晶体中真实的\emph{非中心}（noncentral）多体原子间力。

在共价晶体或金属晶体中，晶格求和（lattice sums）的收敛问题容易安排，因为正离子的长程 Coulomb 场很快被价电子屏蔽掉（参见第 5 章）。但晶格谱的计算会因其他类型的多体力而变得复杂，例如共价键对弯曲（§4.2）、或传导电子气对压缩（§6.11）的抵抗。在非导电晶体中，还必须计入电磁场与极性晶格振动在\emph{极化激元}（polariton）模式中的动力学耦合（§8.3）。

\section{晶格比热}

格波存在的一个重要后果是它们的热激发，它表现为对固体比热的一个贡献。要计算这个贡献，我们其实应当用量子力学的算符来代替经典坐标 $\mathbf{u}_{sl}$ 及其相应的动量。我们在 §2.1 中所做的全部工作，只不过是证明这些坐标可以经过正则变换变到一组新坐标，在新坐标中运动方程就是一群独立简谐振子的运动方程。因此，激发必然属于 Bose–Einstein 类型；一个简正模可以被赋予任意数目的量子，每个量子的能量为 $\hbar\nu_{\mathbf{q}}$，其中 $\nu_{\mathbf{q}}$ 是相应的经典频率。

统计力学理论于是告诉我们，平均而言，第 $q$ 个模中将会有
\begin{equation*}
\bar{n}_{\mathbf{q}} = \frac{1}{e^{\hbar\nu_{\mathbf{q}}/kT}-1}
\tag{2.46}
\end{equation*}
个量子，它们贡献的能量为
\begin{equation*}
\bar{\mathcal{E}}_{\mathbf{q}} = \left(\bar{n}_{\mathbf{q}} + \frac{1}{2}\right)\hbar\nu_{\mathbf{q}},
\tag{2.47}
\end{equation*}
其中计入了零点能（zero-point energy），不过我们暂时把它丢开不算。于是，系统的平均总能量由下式给出：
\begin{equation*}
\bar{\mathcal{E}} = \sum_{\mathbf{q}} \frac{\hbar\nu_{\mathbf{q}}}{e^{\hbar\nu_{\mathbf{q}}/kT}-1},
\tag{2.48}
\end{equation*}
其中求和遍及所有的模（即遍及所有不同的波矢以及所有的偏振）。

由于 $\mathbf{q}$ 矢量在倒格子空间中以密度 $V/8\pi^3$ 分布，我们可以把这个和式写成一个积分；我们还可以通过对温度求微商来计算比热：
\begin{equation*}
C_V = \frac{1}{V}\,\frac{\partial\bar{\mathcal{E}}}{\partial T} = \frac{1}{kT^2}\,\frac{1}{8\pi^3}\sum_{\text{polarizations}} \iiint \frac{(\hbar\nu_{\mathbf{q}})^2\, e^{\hbar\nu_{\mathbf{q}}/kT}}{(e^{\hbar\nu_{\mathbf{q}}/kT}-1)^2}\, d^3q.
\tag{2.49}
\end{equation*}
在形式上，我们可以通过引入\emph{晶格谱}（lattice spectrum）来简化此式。让我们定义 $\mathcal{D}(\nu)$，使得 $\mathcal{D}(\nu)\,d\nu$ 为频率落在 $\nu\to\nu+d\nu$ 范围内的模的分数。对一个每晶胞含 $n$ 个原子的结构，总共有 $3nN$ 个模。于是
\begin{equation*}
C_V = 3nNk \int \frac{(\hbar\nu/kT)^2\, e^{\hbar\nu/kT}}{(e^{\hbar\nu/kT}-1)^2}\,\mathcal{D}(\nu)\, d\nu.
\tag{2.50}
\end{equation*}
但这只不过把问题推回到对 $\mathcal{D}(\nu)$ 的计算上去，而后者可能需要对晶格模的运动方程求出完整的解。采用如下两条假设，我们可以得到一个粗糙的公式。

(i) 只须考虑声学支，并且它们全都具有相同的恒定声速，即
\begin{equation*}
\nu_{\mathbf{q}} = sq.
\tag{2.51}
\end{equation*}

(ii) 限制 $\mathbf{q}$ 允许值范围的布里渊区，可以用倒格子空间中一个同体积的球来代替。

第二个假设意味着，格波存在一个最大波数，即 \emph{Debye 波数}（Debye wave-number）$q_D$。\emph{Debye 球}是布里渊区的一个近似；如果要它恰好容纳 $N$ 个点（$\mathbf{q}$ 空间中点的密度为 $V/8\pi^3$），就必须有
\begin{equation*}
N = \frac{V}{8\pi^3}\,\frac{4}{3}\pi q_D^3,
\tag{2.52}
\end{equation*}
即
\begin{equation*}
q_D = \left(\frac{6\pi^2}{v_c}\right)^{1/3}.
\tag{2.53}
\end{equation*}
在正格子中我们也常常作类似的近似，把 Wigner–Seitz 原胞换成一个半径为 $r_s$ 的 \emph{Wigner–Seitz 球}，其体积与原胞相同：
\begin{equation*}
\frac{4}{3}\pi r_s^3 = v_c.
\tag{2.54}
\end{equation*}
于是
\begin{equation*}
q_D = \left(\frac{9\pi}{2}\right)^{1/3}\frac{1}{r_s}, \qquad \lambda_D = \frac{2\pi}{q_D} \sim 2.6\, r_s.
\tag{2.55}
\end{equation*}
换句话说，声学支的\emph{截止波长}（cut-off wavelength）略大于一个晶胞的平均直径。晶格不能传播波长更短的波。

由这两条假设，晶格谱取如下形式：
\begin{align*}
\mathcal{D}(\nu)\, d\nu &= 4\pi q^2\, dq \,\Big/ \frac{4}{3}\pi q_D^3 \\
&= \frac{3\nu^2}{\nu_D^3}\, d\nu,
\tag{2.56}
\end{align*}
其中 $\nu_D$ 是 \emph{Debye 频率}，$sq_D$。比热公式变成
\begin{align*}
C_V &= 3Nk \int_0^{\nu_D} \frac{(\hbar\nu/kT)^2\, e^{\hbar\nu/kT}}{(e^{\hbar\nu/kT}-1)^2}\,\frac{3\nu^2}{\nu_D^3}\, d\nu \\
&= 3Nk\left(\frac{T}{\Theta}\right)^3 3\int_0^{\Theta/T} \frac{z^4 e^z}{(e^z-1)^2}\, dz,
\tag{2.57}
\end{align*}
其中我们已令 $z=\hbar\nu/kT$，而
\begin{equation*}
k\Theta = \hbar\nu_D
\tag{2.58}
\end{equation*}
定义了 \emph{Debye 温度}（Debye temperature）$\Theta$。

\begin{figure}[H]
\centering
\includegraphics[width=0.72\textwidth]{fig_24.pdf}
\caption*{图 24\quad 比热的 Debye 定律}
\end{figure}

这就是著名的 \emph{Debye 比热公式}。它的结构非常简单。

(i) 当 $T/\Theta$ 大时，积分上限很小，可以把被积函数展开，得到
\begin{equation*}
\int_0^{\Theta/T} \frac{z^4}{z^2}\, dz = \frac{1}{3}\left(\frac{\Theta}{T}\right)^3.
\tag{2.59}
\end{equation*}
于是
\begin{equation*}
C_V = 3Nk,
\tag{2.60}
\end{equation*}
这就是 \emph{Dulong–Petit 定律}。

(ii) 在低温下，积分的上限 $\Theta/T$ 在一切实际场合都可以取作无穷大。这时积分趋于一个常数 $(4\pi^4/15)$，因而
\begin{equation*}
C_V \sim \frac{12\pi^4}{5}\, Nk \left(\frac{T}{\Theta}\right)^3,
\tag{2.61}
\end{equation*}
这就是众所周知的\emph{比热的 $T^3$ 定律}，在低温下成立。

Debye 公式对大多数固体都很有效，Debye 温度已作为固体的一个物理参数被列成表格。它是表征晶格动力学运动的最方便的标度温度：$k\Theta$ 代表在晶格诸模中能够激发的最大量子的能量；$\Theta$ 也应当通过如下关系与晶体中的平均声速相联系：
\begin{equation*}
k\Theta = \hbar q_D\, s.
\tag{2.62}
\end{equation*}

然而，我们在推导这个公式时作了非常粗暴的假设。我们为晶格谱（图 25(a)）假定了一个非常特殊的形式。关系式 $\mathcal{D}(\nu)\propto\nu^2$ 在 $\nu=0$ 附近应当成立，那时材料表现得像一个弹性连续介质，但在
\begin{figure}[H]
\centering
\includegraphics[width=0.72\textwidth]{fig_25.pdf}
\caption*{图 25\quad (a) Debye 谱。(b) 真实的晶格谱}
\end{figure}
$\nu_D$ 处的陡峭截止却并没有根据。精确计算（例如图 25(b)）表明，谱分布在很宽的范围上，带有几个峰，对应于不同偏振的模具有极不相同的速度。高频处往往还有一个相当显著的峰，它源于区边界附近的强色散。$\nu_{\mathbf{q}}$ 随 $\mathbf{q}$ 变化的曲线趋于平坦（参见 §2.2 中的图 22），使得频率相差 $d\nu$ 的诸曲面之间包进了更大的 $\mathbf{q}$ 空间体积。

如果格子是带基元的格子（lattice with a basis），我们还应当计入光学支的贡献。这些模的频率几乎不依赖于波数，因此 \emph{Einstein 模型}（Einstein model）——每个模都具有同一频率——可能更为合适，即
\begin{align*}
C_V &= 3Nk\,\frac{(\hbar\nu_E/kT)^2\, e^{\hbar\nu_E/kT}}{(e^{\hbar\nu_E/kT}-1)^2} \\
&= 3Nk\,\frac{(\Theta_E/T)^2\, e^{\Theta_E/T}}{(e^{\Theta_E/T}-1)^2},
\tag{2.63}
\end{align*}
其中 $\Theta_E$ 是 \emph{Einstein 温度}（Einstein temperature），由下式定义：
\begin{equation*}
k\Theta_E = \hbar\nu_E.
\tag{2.64}
\end{equation*}

然而应当指出，如果把式 (2.57) 中的 $N$ 取作晶体中原子的总数而不是晶胞的数目，我们仍然可以用单独一个 Debye 项来表示复杂晶体的比热。如果我们忽略晶胞内原子彼此之间在位置和质量上的差异，得到的结果正是如此。Debye 公式丝毫不涉及真实的晶体结构，所以这未必是一个坏的近似。实际上，我们正是
利用一个人为扩大的布里渊区——扩展区图式（图 26(a)）——并忽略由此产生的带隙：正如 §2.2 中图 19(c)（原书第 34 页）中那样，这些带隙来自晶胞内各原子之间质量的不同以及动力学相互作用的不同。

\section{晶格谱}

简正模频率分布的实际计算是一个数值计算问题。归根到底，除了对 q 空间中一组密集分布
的 q 值网格求解方程组 (2.13)，再数出有多少个 $\nu_{\mathbf{q}}$ 值落入每个区间
$d\nu$ 之外，别无他法。

但是有一条重要的原理，即 van Hove 定理（van Hove's theorem），它支配着函数
$\mathcal{D}(\nu)$ 及其奇点的性质。这一定理的完整证明超出了本书的范围，但其一般论证
既有意义又重要。

为简单起见，我们考虑谱中的一支。频率落在区间 $d\nu$ 内的模所占的比例等于
\begin{equation*}
\mathcal{D}(\nu)\,d\nu = \frac{v_c}{8\pi^3} \int\!\!\!\int\!\!\!\int d^3 q,
\tag{2.65}
\end{equation*}
其中积分遍及 q 空间中 $\nu \leqslant \nu_{\mathbf{q}} \leqslant \nu + d\nu$
的壳层体积。

我们引入矢量
\begin{equation*}
\mathbf{v}_{\mathbf{q}} \equiv \nabla_{\mathbf{q}}\,\nu_{\mathbf{q}}
\equiv \left( \frac{\partial\nu_q}{\partial q_x},\,
\frac{\partial\nu_q}{\partial q_y},\,
\frac{\partial\nu_q}{\partial q_z} \right),
\tag{2.66}
\end{equation*}
即频率函数在 q 空间中的梯度。这个矢量具有速度的量纲（在 Debye 模型中它就是声速），
并且在所有涉及波的能量传播的问题中都扮演重要角色。事实上，它就是波矢为 $\mathbf{q}$
的波包在晶格这种色散介质中的群速（group velocity）。

我们可以用式 (2.66) 来对式 (2.65) 作一个形式上的简化。令我们在其中积分的 q 空间壳层
内的体积元是一个柱体：其底面是频率面 $\nu_{\mathbf{q}} = \nu$ 上的面元 $dS_\nu$，
高为
\begin{equation*}
dq_\perp = \frac{d\nu}{\left| \partial\nu_{\mathbf{q}}/\partial q \right|}
\tag{2.67}
\end{equation*}
沿曲面法向量出。这个方向正是 $\mathbf{v}_{\mathbf{q}}$ 的方向，于是
\begin{equation*}
\mathcal{D}(\nu)\,d\nu = \frac{1}{8\pi^3 N} \int\!\!\int dS_\nu\, dq_\perp
\tag{2.68}
\end{equation*}

\begin{figure}[H]
\centering
\includegraphics[width=0.72\textwidth]{fig_26.pdf}
\caption*{图 26\quad 声学与光学频率面：(a) 扩展区图式；(b) 分处于两个区中。}
\end{figure}
由此得
\begin{equation*}
\mathcal{D}(\nu) = \frac{1}{8\pi^3 N} \int \frac{dS_\nu}{v_q},
\tag{2.69}
\end{equation*}
其中 $v_q$ 就是矢量 $\mathbf{v}_{\mathbf{q}}$ 的模。

在晶格模谱和电子态谱的理论中我们经常用到这个公式。就其本身而言，它能告诉我们的并不
多，只是关于 $\mathcal{D}(\nu)$ 之奇点的性质。这些奇点必定出现在临界点（critical
point）处，即 $\mathbf{v}_{\mathbf{q}}$ 为零之处，也就是说函数
$\nu_{\mathbf{q}}$ 局部“平坦”之处。

设 $\mathbf{q}_c$ 是一个临界点。由于 $\nu_{\mathbf{q}}$ 是 $\mathbf{q}$
的连续函数，可以在该点附近把它展成 Taylor 级数。线性项因
$\mathbf{v}_{\mathbf{q}} = \mathbf{0}$ 而消失，二次项则可通过主轴变换化为平方和。
于是可以写出
\begin{equation*}
\nu_{\mathbf{q}} = \nu_c + \alpha_1 \xi_1^2 + \alpha_2 \xi_2^2
+ \alpha_3 \xi_3^2 + \ldots,
\tag{2.70}
\end{equation*}
其中 $\boldsymbol{\xi} = \mathbf{q} - \mathbf{q}_c$
是离开临界点的矢距，以局部主轴为参照；系数 $\alpha_1, \alpha_2, \alpha_3$
取决于 $\nu_{\mathbf{q}}$ 对 $\mathbf{q}$ 的局部二阶导数。

例如，设 $\alpha_1, \alpha_2, \alpha_3$ 全为负，则 $\nu$
位于一个局部极大附近。等频面 (2.70) 是椭球面；由初等解析几何可知，
围绕 $\mathbf{q}_c$ 的等值面 $\nu$ 所包围的体积为
\begin{equation*}
\frac{4}{3}\pi\,
\frac{(\nu_c - \nu)^{\frac{3}{2}}}{|\alpha_1 \alpha_2 \alpha_3|^{\frac{1}{2}}},
\tag{2.71}
\end{equation*}
于是由式 (2.65) 出发，对 $\nu$ 求导一次即得
\begin{equation*}
\mathcal{D}(\nu) =
\frac{1}{4\pi^2 N\,|\alpha_1 \alpha_2 \alpha_3|^{\frac{1}{2}}}
(\nu_c - \nu)^{\frac{1}{2}}.
\tag{2.72}
\end{equation*}
此式对 $\nu < \nu_c$ 成立；当 $\nu > \nu_c$ 时，$\mathbf{q}_c$
的邻域对 $\mathcal{D}(\nu)$ 没有贡献。因此，这一奇点并不破坏
$\mathcal{D}(\nu)$ 的连续性，但其斜率 $\partial\mathcal{D}(\nu)/\partial\nu$
不连续，且当 $\nu$ 从下方趋于 $\nu_c$ 时趋于 $-\infty$。

系数 $\alpha_1, \alpha_2, \alpha_3$ 还有其他的可能组合。例如，若一正两负，
我们就得到指标为 1 的鞍点（saddle-point of index 1）。此时，用完全相同的解析几何
方法可以得出，$\nu_c$ 邻域内谱的形状为
\begin{equation*}
\mathcal{D}(\nu) = \left\{
\begin{array}{ll}
C + O(\nu - \nu_c) & (\nu < \nu_c), \\[6pt]
C - \dfrac{1}{4\pi^2 N\,|\alpha_1 \alpha_2 \alpha_3|^{\frac{1}{2}}}
(\nu - \nu_c)^{\frac{1}{2}} + O(\nu - \nu_c) & (\nu > \nu_c),
\end{array}
\right.
\tag{2.73}
\end{equation*}
仍然只是斜率上的一次不连续——一侧带有一条竖直切线——叠加在由区的其他区域产生的光滑
函数 $C + O(\nu - \nu_c)$ 之上。

\begin{figure}[H]
\centering
\includegraphics[width=0.72\textwidth]{fig_27.pdf}
\caption*{图 27\quad 不同类型的 van Hove 奇点。}
\end{figure}
事实上，临界点共有四种类型，如图 27 所示，其中 $S_1$ 和 $S_2$
分别指指标为 1 和指标为 2 的鞍点（后者有两个 $\alpha$
为正、一个为负，其表现犹如反过来的式 (2.73)）。极小的情形显然与极大的情形类似。

这些奇点并不十分严重——但却是必不可少的。van Hove 定理实质上断言：\emph{谱必定至少
各含一个 $S_1$ 型与 $S_2$ 型的临界点，且 $\mathcal{D}(\nu)$
的斜率在上端必趋于 $-\infty$}。这一证明需要若干泛函拓扑方面的定理，超出了本书当前的
范围。不过，这一论证可以在二维情形中得到说明（在二维中，鞍点产生对数奇点，而极值点
产生 $\mathcal{D}(\nu)$ 的有限跃变）。

关键在于，函数 $\nu_{\mathbf{q}}$ 是倒格子空间中 $\mathbf{q}$
的连续周期函数。设在区内某点 $m$ 处它有一个极小，则在倒格子的每个对应点处它都有类似
的极小。同样，如果它在区内某点 $M$ 处取极大值，那么这就是一个临界点，对应于重复区图
式中环绕 $M$ 的山丘的峰顶。\footnote{若极大值位于区的角上，我们也许会觉得积分似乎只应
包含落在区内的那一部分。但来自其他角的贡献会自动补足，凑成像式 (2.72)
那样的完整一项。我们也可以通过平移积分区域使之包含一个完整的峰来保证这一点。倒格子的
\emph{任何}一个晶胞都是可取的积分区域；Brillouin 区只不过是约定俗成的选择。}于是，
整个倒格子空间中将有这样一种峰的分布。

现在考虑一条联结两个极大 $M$、$M$ 的路径。函数 $\nu_{\mathbf{q}}$
必定沿这条路径变化；在路径上的某点 $X$ 处，它取得该路径上的极小值。现在让这条路径在
区内移动，就会生成一系列点 $X, X, X$。这些点构成另一条连续路径；由于 $m$ 与 $m$
是 $\nu_{\mathbf{q}}$ 的绝对极小点，这条路径必定经过这些点。

\begin{figure}[H]
\centering
\includegraphics[width=0.72\textwidth]{fig_28.pdf}
\caption*{图 28\quad 寻找鞍点。}
\end{figure}
但现在我们有了一条从两个极小之间穿过的路径 $mm$。在这条路径上必至少有一个极大
$S$。这个点就是一个鞍点——对诸如 $MM$ 的路径而言它是极小，而对诸如 $mm$
的路径而言它是极大。于是，我们的函数 $\nu_{\mathbf{q}}$ 至少有一个 $S$
型的临界点。

在一般的三维问题中，谱有三个声学支，论证更为复杂。区的原点 $\mathbf{q} = 0$
是所有三支的绝对极小——但函数 $\nu_{\mathbf{q}}$ 在该处不能展成 Taylor
级数，而且速度并不为零。横波的两支也不是彼此分立的，因此必须把
$\nu_{\mathbf{q}}$ 当作 $\mathbf{q}$ 的双值函数来处理。不过，上述形式的定理对
总谱仍然成立。当然，实际的临界点可能远多于这个最少数目。

\section{理想晶体的衍射}

如果晶体中的所有原子都严格处在它们“名下”的格点上，那么与晶体相联系的所有可测物理量
都将是严格的周期函数。例如，一个电子的势能将满足条件
\begin{equation*}
\mathcal{V}(\mathbf{r} + \mathbf{l}) = \mathcal{V}(\mathbf{r}).
\tag{2.74}
\end{equation*}

设有一束快速电子射入晶体。由于被这个势散射，电子将从一个电子态跃迁到另一个电子态。
在微扰论的 Born 近似中，初态 $\Psi_{\mathbf{k}}$ 与末态 $\Psi_{\mathbf{k}'}$
之间的跃迁速率正比于矩阵元
\begin{equation*}
\mathcal{M}_{\mathbf{k}'\mathbf{k}} \equiv \int \Psi_{\mathbf{k}'}^{*}(\mathbf{r})\,
\mathcal{V}(\mathbf{r})\, \Psi_{\mathbf{k}}(\mathbf{r})\, \mathrm{d}\mathbf{r}.
\tag{2.75}
\end{equation*}
的平方。

我们取这些态为平面波态
\begin{equation*}
\Psi_{\mathbf{k}} = e^{i\mathbf{k}\cdot\mathbf{r}}.
\tag{2.76}
\end{equation*}
我们又由式 (1.15) 知道，势可以展成 Fourier 级数
\begin{equation*}
\mathcal{V}(\mathbf{r}) = \sum_{\mathbf{g}} \mathcal{V}_{\mathbf{g}}\,
e^{i\mathbf{g}\cdot\mathbf{r}}.
\tag{2.77}
\end{equation*}
于是
\begin{align*}
\mathcal{M}_{\mathbf{k}'\mathbf{k}}
&= \int e^{-i\mathbf{k}'\cdot\mathbf{r}} \sum_{\mathbf{g}}
\mathcal{V}_{\mathbf{g}}\, e^{i\mathbf{g}\cdot\mathbf{r}}\,
e^{i\mathbf{k}\cdot\mathbf{r}}\, \mathrm{d}\mathbf{r} \\
&= \sum_{\mathbf{g}} \mathcal{V}_{\mathbf{g}} \int
e^{i(\mathbf{k}+\mathbf{g}-\mathbf{k}')\cdot\mathbf{r}}\, \mathrm{d}\mathbf{r} \\
&= \left\{ \begin{array}{ll}
\mathcal{V}_{\mathbf{g}} & \text{若}\quad \mathbf{k}+\mathbf{g}-\mathbf{k}' = 0, \\
0 & \text{其余情形}. \\
\end{array} \right\}
\tag{2.78}
\end{align*}

如果 $\mathbf{k}$ 是固定的——也就是说，入射电子束是单色的且方向高度确定——那么我们
只能在满足
\begin{equation*}
\mathbf{k}' = \mathbf{k} + \mathbf{g},
\tag{2.79}
\end{equation*}
的方向上观察到衍射束，其中 $\mathbf{g}$
是晶体的倒格矢之一。

在这种情形下我们还必须坚持：衍射束的能量与入射能量相同，即两个波矢长度相等。这就给
散射角加上了一个条件。若 $2\theta$ 是 $\mathbf{k}$ 与 $\mathbf{k}'$
之间的夹角，则
\begin{equation*}
|\mathbf{g}| = 2|\mathbf{k}|\sin\theta.
\tag{2.80}
\end{equation*}
为了在倒格子空间中满足这些几何条件，我们作出 \emph{Ewald 球}（Ewald
sphere）（图 29(b)），其半径 $OP$ 等于入射波矢。若把倒格子的原点放在 $P$，则矢量
$\mathbf{g}$ 必定把我们带到 $Q$ 点，$OQ$ 即波矢 $\mathbf{k}'$。因此，只要晶体
相对于入射束的取向恰好使倒格子的一个格点落在球面上，就会发生衍射。

\begin{figure}[H]
\centering
\includegraphics[width=0.72\textwidth]{fig_29.pdf}
\caption*{图 29\quad (a) X 射线衍射；(b) Ewald 作图法。}
\end{figure}
我们可以换一种方式来表述这一点。正如 §1.3 中所指出的，倒格矢的长度正比于它所垂直的
晶面族的面间距 $d$ 的倒数，即
\begin{equation*}
|\mathbf{g}| = \frac{2\pi n}{d},
\tag{2.81}
\end{equation*}
其中 $n$ 是整数（$\mathbf{g}$ 在倒格子三矢组（reciprocal triad）
各轴上分量的最大公约数）。于是，若 $\lambda$ 是入射电子的波长，则式 (2.80) 与
(2.81) 给出
\begin{equation*}
n\lambda = 2d\sin\theta.
\tag{2.82}
\end{equation*}
这就是所谓的 \emph{Bragg 反射定律}（Bragg reflection law）。它很容易通过考察从相继
各晶面上反射的束之间的相位关系推出：要发生相干衍射，附加路径 $ABC$
必须是波长的整数倍（图 30）。

散射矩阵元的大小取决于 $\mathcal{V}_{\mathbf{g}}$。对电子而言，这似乎就是局域势的
各 Fourier 分量——但当我们处理极慢电子时这种看法会导致误解，我们将在 §3.6
再回到这一点。同样的理论也适用于 X 射线的衍射，只是 X 射线响应的是晶体中的局域电子
密度。由于电子密度同样是一个像 (2.74) 那样的周期函数，衍射图样在定性上是相同的。本节
的理论自然进而引出 X 射线晶体结构分析那令人赞叹的复杂性与种种成就。

在中子的情形，散射主要来自原子核：它在空间中的行为像一个 $\delta$
函数型的势，而且随元素不同而有相当大的变化。事实上，由于同一元素的不同同位素可能具有
差别很大的中子
散射截面，并且在晶体中随机混在一起，所以总会有一些\emph{非相干}（incoherent）
散射，它给衍射图样添上一片弥漫的背景。即使晶体中只存在一种同位素，这种情形也可能发生，
因为中子对原子核的自旋取向是敏感的。

\begin{figure}[H]
\centering
\includegraphics[width=0.72\textwidth]{fig_30.pdf}
\caption*{图 30\quad Bragg 衍射。}
\end{figure}
\begin{figure}[H]
\centering
\includegraphics[width=0.72\textwidth]{fig_31.pdf}
\caption*{图 31\quad 反铁磁排列，显示出加倍了的晶胞。}
\end{figure}
中子衍射在分析\emph{反铁磁}（antiferromagnetic）晶体时特别有用。在这类晶体中，原子
具有永久原子磁矩，它们以如下方式极化并排列：一个子晶格上的所有原子磁矩指向一个方向，
另一子晶格上的所有原子磁矩指向相反的方向。我们将在 §§10.6、10.11
中讨论这种系统的磁性质。

这里我们只指出：中子带有磁矩，对每个原子处的局域磁场敏感，因而能够区分磁矩取向不同的
原子。其结果是结构的周期单元变大了，这意味着有效布里渊区的尺寸变小了。新的倒格矢被
引入（可对照图 19），于是衍射图样中出现了新的线，称为磁超格线（magnetic
superlattice lines）。

晶体衍射的这一初等\emph{几何理论}（geometrical
theory）忽略了若干实际上的复杂因素，例如下一节要讨论的热振动的效应。我们还必须考虑
入射束与衍射束在穿过较厚样品时与一层层原子持续相互作用的问题。由此产生若干复杂的效应，
例如\emph{初级消光}（primary extinction）和\emph{摆动效应}（pendulum
effect，即“Pendellösung”），对它们的解释要用到 X 射线衍射的\emph{动力学理论}
（dynamical theory）。其要点是：X
射线（或电子、或中子）在晶体中传播时经受速度色散，在接近相干衍射条件时尤其如此（参见
§3.3）。处在不同色散面上的两束或更多的束可以同时被激发，并在穿过样品时相互干涉。在
\emph{低能电子衍射}（low energy electron diffraction，L.E.E.D.）
的情形，这类效应完全支配着实验现象（参见 §6.9）。

\section{含晶格振动的晶体衍射}

现在考虑更一般的情形：势未必是周期的。设它是各个原子势叠加的结果：
\begin{equation*}
\mathcal{V}(\mathbf{r}) = \sum_{\boldsymbol{l}} \mathcal{V}_a(\mathbf{r} - \mathbf{R}_{\boldsymbol{l}}),
\tag{2.83}
\end{equation*}
其中 $\mathbf{R}_{\boldsymbol{l}}$
是本应处在格点 $\boldsymbol{l}$ 上的那个原子（或离子）的实际位置，而
$\mathcal{V}_a$ 是单个原子的势场。为记号简单起见，我们假定是 Bravais 格子。

于是我们可以像式 (2.75) 那样，计算从态 $\Psi_{\mathbf{k}}$ 散射到态
$\Psi_{\mathbf{k}'}$ 的矩阵元：
\begin{align*}
\mathcal{M}_{\mathbf{k}'\mathbf{k}}
&= \int e^{-i\mathbf{k}'\cdot\mathbf{r}} \sum_{\boldsymbol{l}}
\mathcal{V}_a(\mathbf{r} - \mathbf{R}_{\boldsymbol{l}})\,
e^{i\mathbf{k}\cdot\mathbf{r}}\, \mathrm{d}\mathbf{r} \\
&= \sum_{\boldsymbol{l}} \int e^{i(\mathbf{k}-\mathbf{k}')\cdot\mathbf{r}}\,
\mathcal{V}_a(\mathbf{r} - \mathbf{R}_{\boldsymbol{l}})\, \mathrm{d}\mathbf{r} \\
&= \sum_{\boldsymbol{l}} e^{i(\mathbf{k}-\mathbf{k}')\cdot\mathbf{R}_{\boldsymbol{l}}}
\int e^{i(\mathbf{k}-\mathbf{k}')\cdot(\mathbf{r}-\mathbf{R}_{\boldsymbol{l}})}
\mathcal{V}_a(\mathbf{r} - \mathbf{R}_{\boldsymbol{l}})\, \mathrm{d}\mathbf{r}.
\tag{2.84}
\end{align*}

求和中的每个积分名义上遍及整个空间——至少遍及整个晶体。但我们预计原子势
$\mathcal{V}_a(\mathbf{r})$ 的作用范围非常有限，所以积分结果几乎不会依赖于计量它所
用的中心 $\mathbf{R}_{\boldsymbol{l}}$。我们引入\emph{散射矢量}（scattering vector）
\begin{equation*}
\mathbf{K} = \mathbf{k}' - \mathbf{k};
\tag{2.85}
\end{equation*}
于是矩阵元变为
\begin{equation*}
\mathcal{M}_{\mathbf{k}'\mathbf{k}} = \mathcal{V}_a(\mathbf{K}) \cdot
\frac{1}{N} \sum_{\boldsymbol{l}} e^{-i\mathbf{K}\cdot\mathbf{R}_{\boldsymbol{l}}},
\tag{2.86}
\end{equation*}
其中我们把原子势的 Fourier 变换定义为\footnote{我们引入 $N = 1/v_c$
——宏观单位体积晶体中的晶胞数——以使 $\mathcal{V}_a(\mathbf{K})$
具有能量的量纲。式 (2.84) 中用到的平面波态 (2.76)
显然也是按这个宏观单位体积归一化的。}
\begin{equation*}
\mathcal{V}_a(\mathbf{K}) \equiv \frac{1}{v_c} \int
\mathcal{V}_a(\mathbf{r})\, e^{-i\mathbf{K}\cdot\mathbf{r}}\, \mathrm{d}\mathbf{r}.
\tag{2.87}
\end{equation*}

把式 (2.86) 分解成一个\emph{原子因子}（atomic factor）和一个\emph{结构因子}
（structure factor）是极为有用的。我们先考虑理想晶格的结构因子，即令
\begin{equation*}
\mathbf{R}_{\boldsymbol{l}} = \boldsymbol{l}
= l_1\mathbf{a}_1 + l_2\mathbf{a}_2 + l_3\mathbf{a}_3
\tag{2.88}
\end{equation*}
对所有格点都成立。

把散射矢量写成
\begin{equation*}
\mathbf{K} = K_1\mathbf{b}_1 + K_2\mathbf{b}_2 + K_3\mathbf{b}_3,
\tag{2.89}
\end{equation*}
其中 $\mathbf{b}_1, \mathbf{b}_2, \mathbf{b}_3$
是倒格子三矢组 (1.19)。于是结构因子变为
\begin{align*}
\sum_{\boldsymbol{l}} e^{-i\mathbf{K}\cdot\boldsymbol{l}}
&= \sum_{l_1,l_2,l_3} e^{-i(K_1 l_1 + K_2 l_2 + K_3 l_3)} \\
&= \Bigl( \sum_{0 \leqslant l_1 < L_1} e^{-iK_1 l_1} \Bigr)
   \Bigl( \sum_{0 \leqslant l_2 < L_2} e^{-iK_2 l_2} \Bigr)
   \Bigl( \sum_{0 \leqslant l_3 < L_3} e^{-iK_3 l_3} \Bigr) \\
&= \left( \frac{1 - e^{-iK_1 L_1}}{1 - e^{-iK_1}} \right)
   \left( \frac{1 - e^{-iK_2 L_2}}{1 - e^{-iK_2}} \right)
   \left( \frac{1 - e^{-iK_3 L_3}}{1 - e^{-iK_3}} \right),
\tag{2.90}
\end{align*}
其中，如同式 (1.76) 中那样，$L_1, L_2, L_3$
是沿晶体块体各边的晶胞数。

式 (2.90) 中的每个因子只是 1 的量级，并且随 $\mathbf{K}$ 的变化剧烈起伏。对
$\mathbf{K}$ 的一段取值区间作任何平均都会给出零——除非平均区间包含了所有分母的零点，
即
\begin{equation*}
e^{-iK_1} = e^{-iK_2} = e^{-iK_3} = 1,
\tag{2.91}
\end{equation*}
这只有在 $K_1, K_2, K_3$ 各自都是 $2\pi$
的整数倍时才成立。但根据式 (1.20)，这正是 $\mathbf{K}$
应与倒格矢之一重合的条件。

当 $\mathbf{K}$ 真的等于某个倒格矢时，式 (2.90)
中每一项的分子和分母都变为零。不过这时我们回到原来的求和式，注意到每一项都具有
$\exp(-i\mathbf{g}\cdot\boldsymbol{l})$ 的形式，恰好等于 1，于是求和正好是
$N$。这样我们就得到重要的规则
\begin{equation*}
\frac{1}{N} \sum_{\boldsymbol{l}} e^{-i\mathbf{K}\cdot\boldsymbol{l}}
= \delta_{\mathbf{K}\,\mathbf{g}},
\tag{2.92}
\end{equation*}
其中 $\delta$ 记号在 $\mathbf{K}$ 不是集合 $\{\mathbf{g}\}$
中的矢量时为零。

把式 (2.92) 代回式 (2.86)，便得
\begin{equation*}
\mathcal{M}_{\mathbf{k}'\mathbf{k}}
= \mathcal{V}_a(\mathbf{g})\, \delta_{\mathbf{k}'-\mathbf{k},\,\mathbf{g}},
\tag{2.93}
\end{equation*}
这正是我们先前的结果 (2.78)。

但现在假设晶体的晶格模被激发了。每个原子都偏离其理想格点：
\begin{align*}
\mathbf{R}_{\boldsymbol{l}} &= \boldsymbol{l} + \mathbf{u}_{\boldsymbol{l}} \\
&= \boldsymbol{l} + \sum_{\mathbf{q}>}
\left( \mathbf{U}_{\mathbf{q}}\, e^{i\mathbf{q}\cdot\boldsymbol{l}}
+ \mathbf{U}_{\mathbf{q}}^{*}\, e^{-i\mathbf{q}\cdot\boldsymbol{l}} \right),
\tag{2.94}
\end{align*}
其中 $\mathbf{U}_{\mathbf{q}}$ 如式 (2.8) 与 (2.10)
中那样，是波矢为 $\mathbf{q}$ 的模的矢量振幅。求和只遍及 $\mathbf{q}$
取值范围的一半，这样我们就可以把复共轭项 $\mathbf{U}_{\mathbf{q}}^{*} =
\mathbf{U}_{-\mathbf{q}}$ 明写出来，使位移为实。

把式 (2.94) 代入式 (2.86) 的结构因子中：
\begin{align*}
\frac{1}{N} \sum_{\boldsymbol{l}} e^{-i\mathbf{K}\cdot\mathbf{R}_{\boldsymbol{l}}}
&= \frac{1}{N} \sum_{\boldsymbol{l}} \exp\bigl[ -i\mathbf{K}\cdot
\bigl\{ \boldsymbol{l} + \sum_{\mathbf{q}>}
(\mathbf{U}_{\mathbf{q}}\, e^{i\mathbf{q}\cdot\boldsymbol{l}}
+ \mathbf{U}_{\mathbf{q}}^{*}\, e^{-i\mathbf{q}\cdot\boldsymbol{l}}) \bigr\} \bigr] \\
&= \frac{1}{N} \sum_{\boldsymbol{l}} \bigl[ e^{-i\mathbf{K}\cdot\boldsymbol{l}}
\prod_{\mathbf{q}>} \exp\bigl\{ -i\mathbf{K}\cdot
(\mathbf{U}_{\mathbf{q}}\, e^{i\mathbf{q}\cdot\boldsymbol{l}}
+ \mathbf{U}_{\mathbf{q}}^{*}\, e^{-i\mathbf{q}\cdot\boldsymbol{l}}) \bigr\}
\bigr].
\tag{2.95}
\end{align*}

这是严格的。让我们把宗量正比于 $\mathbf{U}_{\mathbf{q}}$
的那个指数函数展开——假定它是小位移：
\begin{align*}
&\exp \bigl\{ -i\mathbf{K}\cdot
(\mathbf{U}_{\mathbf{q}}\, e^{i\mathbf{q}\cdot\boldsymbol{l}}
+ \mathbf{U}_{\mathbf{q}}^{*}\, e^{-i\mathbf{q}\cdot\boldsymbol{l}}) \bigr\} \\
&\qquad = 1 - i\mathbf{K}\cdot
(\mathbf{U}_{\mathbf{q}}\, e^{i\mathbf{q}\cdot\boldsymbol{l}}
+ \mathbf{U}_{\mathbf{q}}^{*}\, e^{-i\mathbf{q}\cdot\boldsymbol{l}})
- |\mathbf{K}\cdot\mathbf{U}_{\mathbf{q}}|^2 \ldots
\tag{2.96}
\end{align*}

代入式 (2.95) 后显然可见，把乘积展开时，会出现一个形如 (2.90)
的项，给出理想晶体结构的衍射图样；
其后是所有与 $\mathbf{K}\cdot\mathbf{U}_{\mathbf{q}}$
呈线性关系的项之和，即形如
\begin{align*}
\frac{1}{N} \sum_{\boldsymbol{l}} e^{-i\mathbf{K}\cdot\boldsymbol{l}}
\sum_{\mathbf{q}} \left( -i\mathbf{K}\cdot\mathbf{U}_{\mathbf{q}} \right)
e^{i\mathbf{q}\cdot\boldsymbol{l}}
&= \frac{1}{N} \sum_{\mathbf{q}} \left( -i\mathbf{K}\cdot\mathbf{U}_{\mathbf{q}} \right)
\sum_{\boldsymbol{l}} e^{i(\mathbf{q}-\mathbf{K})\cdot\boldsymbol{l}} \\
&= \sum_{\mathbf{q}} \left( -i\mathbf{K}\cdot\mathbf{U}_{\mathbf{q}} \right)
\delta_{\mathbf{K}-\mathbf{q},\,\mathbf{g}}
\tag{2.97}
\end{align*}
的项——最后一步用了式 (2.92)。正如我们将在 §2.9
中看到的，式 (2.96) 中的模平方项也很重要，因为它的符号不变。

为了理解这一结果，让我们考虑两种情形。在图 32(a) 中，$\mathbf{K}$
位于布里渊区之内。也就是说，我们在这样的态的区域中寻找被散射的电子（或 X
射线、中子）：其中 $\mathbf{k}'$ 与初态 $\mathbf{k}$ 相差不超过半个倒格矢。按惯例，波矢
$\mathbf{q}$ 被限制在一个布里渊区内，于是式 (2.97)
中剩下的唯一一项就是 $\mathbf{q} = \mathbf{K}$ 的那一项。换言之，这一跃迁的矩阵元为
\begin{equation*}
\mathcal{M}_{\mathbf{k}'\mathbf{k}} = -i\{(\mathbf{k}'-\mathbf{k})\cdot
\mathbf{U}_{\mathbf{q}}\}\, \mathcal{V}_a(\mathbf{k}'-\mathbf{k}),
\tag{2.98}
\end{equation*}
其中
\begin{equation*}
\mathbf{q} = \mathbf{k}' - \mathbf{k}.
\tag{2.99}
\end{equation*}

\begin{figure}[H]
\centering
\includegraphics[width=0.72\textwidth]{fig_32.pdf}
\caption*{图 32\quad (a) 正常过程；(b) 倒逆过程（Umklapp 过程）。}
\end{figure}

如果 $\mathbf{K}$ 并\emph{不}落在第一布里渊区内（图 32(b)），式 (2.97)
的求和式中仍然会剩下一项。也就是说，仍然存在一个 $\mathbf{q}$
值满足 $\mathbf{K} - \mathbf{q} = \mathbf{g}$，即满足
\begin{equation*}
\mathbf{q} = \mathbf{k}' - \mathbf{k} - \mathbf{g}.
\tag{2.100}
\end{equation*}

布里渊区作为倒格子的一个晶胞，可以重复铺排而恰好盖满矢量 $\mathbf{K}$ 的整个空间；$\mathbf{K}$ 的任何一个值都可以唯一地由一个 $\mathbf{g}$ 值——倒格子的一个平移矢——和一个 $\mathbf{q}$ 值构造出来，$\mathbf{q}$ 是一个被限制在该格子一个晶胞之内的矢量。用 §1.5 的语言来说，$\mathbf{q}$ 就是 $\mathbf{K}$ \emph{约化}（reduced）到第一布里渊区后的值。

因此，在这两种情形下式 (2.98) 都成立，其中 $\mathbf{q}$ 由 (2.99) 或 (2.100) 给出，视何者适用而定；如果我们把 $\mathbf{g}=0$ 也包括进倒格矢的完全集合之中，则关系式 (2.100) 恒成立。无论我们选取哪个方向，总会有一些散射，只不过其振幅将取决于碰巧具有满足这些条件的正确波矢的那些晶格振动的振幅。

\section{声子}

一个散射过程能否发生，并不单由矩阵元异于零所决定；还有能量条件需要满足。在理想晶格的情形，整个系统是静态的，因此初态和末态的电子（X 射线、中子）具有相同的能量。

的确，晶格处于不断的运动之中。(2.94) 中的矢量 $\mathbf{U}_{\mathbf{q}}$ 应当含有形如 $\exp(i\nu_{\mathbf{q}}t)$ 的时间因子，其中 $\nu_{\mathbf{q}}$ 是波矢为 $\mathbf{q}$ 的简正模的频率。实际上，$\mathbf{U}_{\mathbf{q}}$ 是三个这类矢量之和，它们的振幅、方向和频率各不相同，对应于该波矢的三个不同的声学模。但我们一次只取其中一个，并且记住，入射束与散射束的态 $\Psi_{\mathbf{k}}$ 和 $\Psi_{\mathbf{k}'}$ 必须具有能量 $\mathcal{E}(\mathbf{k})$ 和 $\mathcal{E}(\mathbf{k}')$，从而带有形如 $\exp\{i\mathcal{E}(\mathbf{k})t/\hbar\}$ 与 $\exp\{i\mathcal{E}(\mathbf{k}')t/\hbar\}$ 的时间因子。当我们对时间求平均时，将会遇到如下形式的积分
\begin{equation*}
\int \exp[i\{\mathcal{E}(\mathbf{k})-\mathcal{E}(\mathbf{k}')+\hbar\nu_{\mathbf{q}}\}t/\hbar]\,dt,
\end{equation*}
除非
\begin{equation*}
\mathcal{E}(\mathbf{k})-\mathcal{E}(\mathbf{k}')+\hbar\nu_{\mathbf{q}}=0.
\tag{2.101}
\end{equation*}
否则它们恒为零。换言之，衍射是\emph{非弹性}的；束流获得了与之相互作用的那个晶格模的一个能量量子。

其实这还不完全，因为晶格振动中还会出现共轭的时间依赖 $\exp(-i\nu_{\mathbf{q}}t)$，从而给出电子损失能量 $\hbar\nu_{\mathbf{q}}$ 的衍射。在这两种情形下，我们上面勾画的论证都可以借助于标准的时间依赖微扰理论而变得完全严格。

条件 (2.100) 与 (2.101) 可以给出一种非常直观的解释。我们说，入射电子（或 X 射线、中子）与晶格相互作用而消灭（或产生）了一个\emph{声子}（phonon），其波矢为 $\mathbf{q}$，能量为 $\hbar\nu_{\mathbf{q}}$。我们用图 33$(a)$ 或 $(b)$ 那样的图形来表示这些过程。条件 (2.101) 就是该过程中的能量守恒规则。

如果我们考察对第一区中 $(\mathbf{k}'-\mathbf{k})$ 的一切值都成立的条件 (2.99)，就会发现它看上去像一条动量守恒定律。乘以 $\hbar$，便得
\begin{equation*}
\hbar\mathbf{k}'=\hbar\mathbf{k}+\hbar\mathbf{q}.
\tag{2.102}
\end{equation*}
电子的动量吸收了被消灭的那个声子的\emph{晶体动量}（crystal momentum）$\hbar\mathbf{q}$。这正是为“声子”这个名称辩护的论证——声子是一种具有“类粒子”性质的量子化声激发，与“光子”相类比。

\begin{figure}[H]
\centering
\includegraphics[width=0.72\textwidth]{fig_33.pdf}
\caption*{图 33\quad 电子散射过程：(a) 声子吸收；(b) 声子发射。}
\end{figure}

但在一般情形下，这条“动量守恒规则”并不成立；正如 (2.100) 所表明的，除声子的动量之外，电子还可以获得或失去动量 $\hbar\mathbf{g}$。我们称这样的过程为倒逆过程（Umklapp process），或 U 过程（U-process）。特殊情形 (2.102) 则可称为正常过程（Normal process），或 N 过程（N-process）。

这个结果其实并不令人惊讶。额外的动量 $\hbar\mathbf{g}$ 只不过是传给了作为整体的晶体。在构造晶格振动态时，我们忽略了理想晶格质心的运动——晶格振动正是相对于它来度量的。可以查验，只要把与这个自由度相联系的运动包括进来，基本的动量守恒定律总体上仍然得到满足。一个倒逆过程可以看成是在产生（或消灭）一个声子的同时发生了一次 Bragg 反射；后一过程中的动量显然传给了作为整体的晶体。

非弹性衍射现象是研究晶体晶格动力学的一种极有价值的工具。沿特定方向衍射的束流与具有确定波矢 $\mathbf{q}$ 的晶格模相联系。我们可以考察衍射粒子能量的变化，从而测出 $\hbar\nu_{\mathbf{q}}$。通过朝不同方向观察，并把晶体转到不同的取向，就可以把整个函数 $\nu_{\mathbf{q}}$ 测绘出来。当然，它有若干个不同偏振的分支，但通过系统分析可以把它们分开。中子对局部磁矩的敏感性（§2.6）意味着，对极化中子束被磁振子（magnons）非弹性衍射的情形也有一个类似的理论（见 §10.10），从而给出关于磁振子色散性质的类似信息。

然而，这个实验只有用“热”中子才切实可行；这种中子在能量约为 0.1 eV 时，其波长与晶格间距同数量级。由声子能量引起的频移——其数量级为 $k\Theta$ 或更小，也许 0.01 eV——很容易观察到。对于电子和 X 射线，束流的能量必须高得多——几十或几百电子伏——因此衍射过程中微小的能量变化无法被探测到。

晶体动量的选择定则 (2.102) 依赖于晶格的长程序。当这种有序不存在时，比如在液体或玻璃中，我们会观察到什么？弹性衍射将依赖于矩阵元 (2.86) 的平方，后者正比于结构因子的平方：
\begin{align*}
S(\mathbf{K}) &\equiv \left|\frac{1}{N}\sum_{\mathbf{l}} e^{-i\mathbf{K}\cdot\mathbf{R}_{\mathbf{l}}}\right|^{2}\\
&= \frac{1}{N^{2}}\sum_{\mathbf{l}\mathbf{l}'} e^{-i\mathbf{K}\cdot(\mathbf{R}_{\mathbf{l}}-\mathbf{R}_{\mathbf{l}'})}\\
&= \frac{1}{N}\int P(\mathbf{R})\,e^{-i\mathbf{K}\cdot\mathbf{R}}\,\mathrm{d}\mathbf{R}.
\tag{2.103}
\end{align*}
衍射图样依赖于对关联函数（pair correlation function）$P(\mathbf{R})$——即在相距 $\mathbf{R}$ 处找到两个原子的概率——的 Fourier 变换。在液体中，这正是我们熟悉的体系的径向分布函数（radial distribution function）。

对这个论证作一个简单的推广即可证明，非弹性衍射的振幅 $S(\mathbf{K},\nu)$ 必定是含时关联函数（time-dependent correlation function）$P(\mathbf{R},t)$ 的双重 Fourier 变换，$P(\mathbf{R},t)$ 度量的是这样的概率：已知在时刻 0 有一个原子（可能就是同一个）位于原点 0，在时刻 $t$ 于点 $\mathbf{R}$ 处找到一个原子的概率。例如，中子衍射可以测量波矢为 $\mathbf{K}$ 的密度涨落的频谱——这个物理量恰好只有在相当完善的晶格中才相当尖锐。类似地，极化中子会被自旋密度（spin density）涨落所衍射，这些涨落由类似的关联函数或它们在空间或动量、时间或能量上的 Fourier 变换所定义。

\section{德拜–沃勒因子}

(2.97) 中那些对应于含一个声子的散射过程的项，并不是结构因子 (2.95) 中可能出现的一切。容易看出，类似 (2.96) 的那些因子的乘积，以及指数展开中更高阶的项，会产生包含形如 $\mathbf{U}_{\mathbf{q}}\exp(i\mathbf{q}\cdot\mathbf{l})$ 的因子的各种乘积的贡献。对每一个这样的因子，我们都称相应的声子被产生或消灭了，因此这些项属于多声子过程（multiphonon processes）。

总的来说，随着阶数升高，这些过程的速率迅速下降，对非弹性衍射背景的贡献并不很大。然而，有一类重要的项来自展开式 (2.96) 中的平方项 $-|\mathbf{K}\cdot\mathbf{U}_{\mathbf{q}}|^{2}$，它们平均之后并不给出小的贡献。考察 (2.95) 与 (2.96) 即可发现，无论是弹性还是非弹性散射，矩阵元的平方都应当乘以
\begin{equation*}
e^{-2W}=\prod_{\mathbf{q}}\left\{1-|\mathbf{K}\cdot\mathbf{U}_{\mathbf{q}}|^{2}\right\},
\tag{2.104}
\end{equation*}
现在 $\mathbf{q}$ 则重新取遍整个布里渊区。

这个量称为 Debye–Waller 因子（Debye–Waller factor）。把它写成这种形式，是因为我们可以利用一条标准的代数定理
\begin{equation*}
\lim_{N\to\infty}\prod_{n=1}^{N}\left(1-\frac{1}{N}a_{n}\right)=\exp\left\{-\lim_{N\to\infty}\frac{1}{N}\sum_{n=1}^{N}a_{n}\right\},
\tag{2.105}
\end{equation*}
把乘积变换成求和。于是
\begin{equation*}
e^{-2W}=\exp\left\{-\sum_{\mathbf{q}}\frac{1}{2}\,|\mathbf{K}\cdot\mathbf{U}_{\mathbf{q}}|^{2}\right\},
\tag{2.106}
\end{equation*}
即
\begin{equation*}
W=\frac{1}{2}\sum_{\mathbf{q}}|\mathbf{K}\cdot\mathbf{U}_{\mathbf{q}}|^{2}.
\tag{2.107}
\end{equation*}
要计算这个和，我们需要知道第 $\mathbf{q}$ 个晶格模的振幅 $\mathbf{U}_{\mathbf{q}}$。它将是温度的函数。我们知道，该模中的平均能量由 (2.47) 给出：
\begin{equation*}
\bar{\mathcal{E}}_{\mathbf{q}}=\left(\bar{n}_{\mathbf{q}}+\frac{1}{2}\right)\hbar\nu_{\mathbf{q}},
\end{equation*}
其中 $\bar{n}_{\mathbf{q}}$ 是该模中声子的平均数，由 Bose–Einstein 公式 (2.46) 给出。

在经典处理中，我们可以把每个简谐振子模的能量算成其动能与势能之和。众所周知，这两者相等——于是由 (2.1) 和 (2.8) 得
\begin{align*}
\bar{\mathcal{E}}&=\sum_{s\mathbf{l}}M_{s}\,|\dot{\mathbf{u}}_{s\mathbf{l}}|^{2}\\
&=\sum_{s\mathbf{q}}NM_{s}\,|\dot{\mathbf{U}}_{s\mathbf{q}}|^{2}\\
&=\sum_{s\mathbf{q}}NM_{s}\,\nu_{\mathbf{q}}^{2}|\mathbf{U}_{s\mathbf{q}}|^{2}.
\tag{2.108}
\end{align*}
如果每个晶胞只有一个原子，质量为 $M$，则
\begin{align*}
|\mathbf{U}_{\mathbf{q}}|^{2}&=\bar{\mathcal{E}}_{\mathbf{q}}/NM\nu_{\mathbf{q}}^{2}\\
&=(\bar{n}_{\mathbf{q}}+\tfrac{1}{2})\hbar/NM\nu_{\mathbf{q}}.
\tag{2.109}
\end{align*}
我们知道 $\mathbf{U}_{\mathbf{q}}$ 在晶格谱每一支上的偏振，所以原则上能够精确计算出 Debye–Waller 因子。

为了看出它应当有怎样的行为，让我们假定一个三个模都具有相同速度的 Debye 模型。对任一偏振，平均而言我们应得
\begin{equation*}
|\mathbf{K}\cdot\mathbf{U}_{\mathbf{q}}|^{2}=\tfrac{1}{3}K^{2}|\mathbf{U}_{\mathbf{q}}|^{2},
\tag{2.110}
\end{equation*}
但若计入三个不同的偏振，因子 $\frac{1}{3}$ 便消去了。\footnote{有一条有用的规则：在这个模型中，三个模对每个 $\mathbf{q}$ 值都是简并的。因此可以随意选取偏振矢量，只要它们相互正交。于是可以取一个模的 $\mathbf{U}_{\mathbf{q}}$ 平行于 $\mathbf{K}$，另外两个垂直于 $\mathbf{K}$。这就给出了我们所需要的结果。}

由 (2.46)、(2.56) 和 (2.109) 得
\begin{align*}
W&=\frac{1}{2}K^{2}\frac{\hbar}{NM}\frac{1}{8\pi^{3}}\iiint\frac{\bar{n}_{\mathbf{q}}+\frac{1}{2}}{\nu_{\mathbf{q}}}\,d^{3}q\\
&=\frac{1}{2}\frac{\hbar^{2}K^{2}}{M}\int_{0}^{\nu_{D}}\left\{\frac{1}{e^{\hbar\nu/kT}-1}+\frac{1}{2}\right\}\frac{3\nu^{2}}{\hbar\nu\,\nu_{D}^{3}}\,d\nu\\
&=\frac{3}{2}\frac{\hbar^{2}K^{2}T^{2}}{Mk\Theta^{3}}\int_{0}^{\Theta/T}\left\{\frac{1}{e^{z}-1}+\frac{1}{2}\right\}z\,dz,\quad\text{正如式 (2.57) 中那样。}
\tag{2.111}
\end{align*}
高温下积分的上限很小，被积函数中的指数因子可以按 $z$ 的幂展开。结果是
\begin{equation*}
W\to\frac{3}{2}\frac{\hbar^{2}K^{2}T}{Mk\Theta^{2}}.
\tag{2.112}
\end{equation*}
于是，正比于矩阵元 $\mathcal{M}_{\mathbf{k}'\mathbf{k}}$ 的平方的 X 射线衍射图样，其强度将被如下因子所削减：
\begin{equation*}
e^{-2W}\sim\exp(-3\hbar^{2}K^{2}T/Mk\Theta^{2}).
\tag{2.113}
\end{equation*}
这个因子相当强地依赖于温度，也依赖于散射矢量的大小。如果我们对每个模的平均能量采用经典统计，即 $\bar{\mathcal{E}}_{\mathbf{q}}=kT$，也会得到这个结果。

在低于 Debye 温度的温度下，公式要更复杂一些，但我们注意到，在极低温下 $W$ 将趋于一个常数。这源于积分中的 $\frac{1}{2}$ 项——一项来自晶格的零点运动（zero-point motion）的贡献。这并不是一个可以忽略的效应：
\begin{equation*}
W\to\frac{3}{8}\frac{\hbar^{2}K^{2}}{Mk\Theta}\quad\text{当}\ T\to 0\ \text{时},
\tag{2.114}
\end{equation*}
它是 $W$ 在 $T\sim\Theta$ 处取值的 $\frac{1}{4}$。零点\emph{能}在物理上也许无足轻重，但与之相联系的\emph{运动}却可以被直接观察到。

非常类似的论证可以解释 Mössbauer 效应（Mössbauer effect）。原子核 γ 发射的谱中可能包含一条极窄的谱线，看不出它被原子在晶体中绕其平衡位置的热运动所展宽。这其实是一个简单的经典现象。以匀速 $v$ 运动的源所发射的频率为 $\omega$ 的辐射，会经受大小为 $\omega v/c$ 的 Doppler 频移。但设想这个源由一个束缚原子携带，该原子以频率 $\nu_{\mathbf{q}}$、振幅 $\mathbf{U}_{\mathbf{q}}$ 往复振动。由于速度现在随时间变化，我们观察到的发射辐射是调频的。正如调频（FM）广播的初等理论那样，谱现在分裂成位于 $\omega$、$\omega\pm\nu$、$\omega\pm2\nu$ 等频率处的若干条线。对于小的振动振幅，未发生改变的那条谱线在谱中所占的份额近似为 $1-\frac{1}{2}(\mathbf{k}\cdot\mathbf{U}_{\mathbf{q}})^{2}$，其中 $\mathbf{k}$ 是 γ 光子的波矢。像 (2.104–7) 那样把所有可能振动模的贡献合并起来，我们发现这条谱线的相对强度是一个有限量，由一个正像 Debye–Waller 因子一样的公式给出。

用固体理论的标准语言，我们可以说，γ 射线的发射传给原子的冲量有一部分被声子的产生所吸收，声子的能量则表现为谱线的展宽。但有有限比例的动量传给了作为整体的晶体，不伴随声子的产生（或吸收），因而只有无穷小的能量损失（或获得）。正如我们对 X 射线衍射的分析那样，Debye–Waller 因子度量的正是这类弹性即无反冲（recoilless）事件所占的比例。这条极锐的谱线当然极其有用：可以用来探测分子与固体中的磁场，检验相对论理论，等等。

(2.111) 的推导还顺带给了我们另一个有趣的结果。由 (2.1) 和 (2.8) 显然可见，在 Debye 温度以上，每个原子绕其格点振动的均方振幅由下式给出：
\begin{align*}
\frac{1}{N}\sum_{\mathbf{l}}|\mathbf{u}_{\mathbf{l}}|^{2}&=\sum_{\mathbf{q}}|\mathbf{U}_{\mathbf{q}}|^{2}\\
&\approx\frac{9\hbar^{2}T}{Mk\Theta^{2}}
\tag{2.115}
\end{align*}
于是在温度 $T$ 下，每个原子偏离其平衡位置的均方根位移将是，比如说，晶胞平均半径 $r_{s}$ 的一个分数 $x$，其中
\begin{equation*}
x=\sqrt{\frac{9\hbar^{2}T}{Mk\Theta^{2}r_{s}^{2}}}.
\tag{2.116}
\end{equation*}
Lindemann 熔化公式（Lindemann melting formula）正是基于这个观念。它认为，当 $x$ 达到某个标准值 $x_{m}$ 时，固体就必须熔化。于是，熔化温度 $T_{m}$ 便与固体的其他原子常数通过下式相联系：
\begin{equation*}
T_{m}=\frac{x_{m}^{2}}{9\hbar^{2}}\,Mk\Theta^{2}r_{s}^{2}.
\tag{2.117}
\end{equation*}
看来在大多数固体中 $x_{m}$ 处于 0.2–0.25 的范围内。

这条规则可以用来在已知 $T_{m}$ 的情形下近似地估计 Debye 温度。它也为估计像 $W$ 这类依赖于原子振动振幅的量的数值提供了一条便捷的捷径。例如，我们可以证明，(2.112) 可以写成
\begin{equation*}
W\approx x_{m}^{2}\,\frac{T}{T_{m}}\,\frac{K^{2}}{q_{D}^{2}},
\tag{2.118}
\end{equation*}
其中 $q_{D}$ 是 Debye 波数 (2.53)。

\section{非谐性与热膨胀}

在本章开始时，我们把晶体的势能展开成了晶格位移的 Taylor 级数。但这个级数——式 (2.2)——在二阶项处被截断了。还会有更进一步的项，例如
\begin{equation*}
\mathcal{V}^{(3)}=\frac{1}{3!}\sum_{\substack{ss's''\\ ll'l''\\ jj'j''}}u_{s\mathbf{l}}^{j}\,u_{s'\mathbf{l}'}^{j'}\,u_{s''\mathbf{l}''}^{j''}\left[\frac{\partial^{3}\mathcal{V}}{\partial u_{s\mathbf{l}}^{j}\,\partial u_{s'\mathbf{l}'}^{j'}\,\partial u_{s''\mathbf{l}''}^{j''}}\right]_{0},
\tag{2.119}
\end{equation*}
如此等等。这些系数的实际计算是一个非常复杂的问题，因为它们既涉及几何因子，又涉及原子间势的三阶导数。

有若干重要的物理现象与非谐项（anharmonic terms）相联系。其中最熟知的是热膨胀（thermal expansion）。要从 (2.119) 这类表达式直接导出它并不容易，但一般的物理图像却相当简单。随着温度升高，晶格振动的振幅增大，因而位移 $\mathbf{u}_{s\mathbf{l}}$ 等的均方根平均值增大。非谐项对晶体的自由能有贡献，于是自由能不再必然在围绕所假定的“平衡”位形（其中每个 $\mathbf{u}_{s\mathbf{l}}$ 均为零）的振动下取极小。整个晶体于是会膨胀（或收缩），直至找到总自由能取极小的那个体积。

由于缺乏关于非谐项的充分知识而无法作“第一性原理”的计算，我们可以唯象地假定晶格诸模的频率是体积的函数。为简单起见，设体积变化 $\Delta V$ 使每个模的频率产生相同的相对变化：
\begin{equation*}
\frac{\Delta\nu}{\nu}=-\gamma\,\frac{\Delta V}{V}.
\tag{2.120}
\end{equation*}
于是，晶体总自由能作为体积的函数便可写成
\begin{equation*}
F=\frac{1}{2}\,\frac{1}{\kappa}\left(\frac{\Delta V}{V}\right)^{2}+kT\sum_{\mathbf{q}}\ln\left(2\sinh\frac{\hbar\nu_{\mathbf{q}}}{2kT}\right).
\tag{2.121}
\end{equation*}
第一项是与作为弹性连续介质的固体的压缩率（compressibility）$\kappa$ 相联系的势能。第二项是诸晶格模自由能之和，由 Bose–Einstein 振子的常规统计力学给出。

对体积求微商，并利用 (2.120)，我们得到自由能取极小的条件：
\begin{align*}
\frac{1}{\kappa}\left(\frac{\Delta V}{V}\right)&=\sum_{\mathbf{q}}\gamma\hbar\nu_{\mathbf{q}}\,\tfrac{1}{2}\coth\frac{\hbar\nu_{\mathbf{q}}}{2kT}\\
&=\gamma\bar{\mathcal{E}}(T),
\tag{2.122}
\end{align*}
其中 $\bar{\mathcal{E}}$ 是温度 $T$ 下诸晶格模中的能量，如 (2.48) 所示。这样我们便得到了 Grüneisen 公式（Grüneisen formula）：温度 $T$ 下的体膨胀（dilatation）正比于平均热能量密度，即
\begin{equation*}
\frac{\Delta V}{V}=\kappa\gamma\bar{\mathcal{E}}(T).
\tag{2.123}
\end{equation*}
热膨胀系数（thermal expansion coefficient）作为体膨胀对温度的导数，正比于比热 $C_{V}$。

把这个公式与实验比较，可以定出 Grüneisen 常量（Grüneisen constant）$\gamma$，其值通常约为 2。它是刻画非谐性效应的一个方便的无量纲参量。在 Debye 模型中，我们可以把 (2.120) 写成
\begin{equation*}
\gamma=-\frac{\partial\ln\Theta}{\partial\ln V}
\tag{2.124}
\end{equation*}
的形式，它表明了体积对 Debye 温度的影响。说实话，作为热膨胀的一个模型，这实在过于简化了。体膨胀对不同模的影响方式并不相同；纵模的 $\gamma$ 值通常远大于横模的 $\gamma$ 值，因此在 (2.122) 中必须把它们作为单独的贡献分别计入。

非谐性的另一个效应表现在弹性常数上，弹性常数随体积和温度而变化。这些是复杂的现象，没有初等理论可循；它们依赖于若干个不同的参量。

\section{声子–声子相互作用}

从形式上看，非谐项最重要的效应，是它破坏了各个晶格模的动力学独立性。例如，假设我们把代换 (2.8) 与 (2.10) 代入 (2.119)。对于 Bravais 格子，可以写出
\begin{align*}
\mathcal{V}^{(3)} &= \sum_{ll'l''}\mathsf{A}_{ll'l''}\colon \mathbf{u}_l\,\mathbf{u}_{l'}\,\mathbf{u}_{l''}\\
&= \sum_{ll'l''}\sum_{\mathbf{q}\mathbf{q}'\mathbf{q}''}\mathsf{A}_{ll'l''}\colon \mathbf{U}_{\mathbf{q}}\mathbf{U}_{\mathbf{q}'}\mathbf{U}_{\mathbf{q}''}\exp\bigl\{i(\mathbf{q}\cdot\mathbf{l}+\mathbf{q}'\cdot\mathbf{l}'+\mathbf{q}''\cdot\mathbf{l}'')\bigr\},
\tag{2.125}
\end{align*}
这里我们对三阶系数张量、以及它与各个晶格位移矢量的标积，都使用了非常紧凑的记号。

这一项本应计入运动方程之中。我们可以像 (2.6) 中那样应用平移不变性原理，来说明这类项的性质。(2.125) 中张量的系数只应依赖于格点 $\mathbf{l}$、$\mathbf{l}'$ 和 $\mathbf{l}''$ 的相对位置。于是
\begin{equation*}
\mathsf{A}_{ll'l''}=\mathsf{A}(\mathbf{h}',\mathbf{h}'')\quad\text{其中}\quad \mathbf{h}'=\mathbf{l}'-\mathbf{l},\ \mathbf{h}''=\mathbf{l}''-\mathbf{l}:
\tag{2.126}
\end{equation*}
把它代入 (2.125)，便得到
\begin{align*}
\mathcal{V}^{(3)} &= \sum_{l\mathbf{h}'\mathbf{h}''}\sum_{\mathbf{q}\mathbf{q}'\mathbf{q}''}\mathsf{A}(\mathbf{h}',\mathbf{h}'')\colon \mathbf{U}_{\mathbf{q}}\mathbf{U}_{\mathbf{q}'}\mathbf{U}_{\mathbf{q}''}\,e^{i(\mathbf{q}+\mathbf{q}'+\mathbf{q}'')\cdot\mathbf{l}}\,e^{i\mathbf{q}'\cdot\mathbf{h}'}e^{i\mathbf{q}''\cdot\mathbf{h}''}\\
&= \sum_{\mathbf{q}\mathbf{q}'\mathbf{q}''}\Bigl\{\sum_{\mathbf{h}'\mathbf{h}''}\mathsf{A}(\mathbf{h}',\mathbf{h}'')\,e^{i\mathbf{q}'\cdot\mathbf{h}'}e^{i\mathbf{q}''\cdot\mathbf{h}''}\Bigr\}\colon \mathbf{U}_{\mathbf{q}}\mathbf{U}_{\mathbf{q}'}\mathbf{U}_{\mathbf{q}''}\sum_{l}e^{i(\mathbf{q}+\mathbf{q}'+\mathbf{q}'')\cdot\mathbf{l}}\\
&= \sum_{\mathbf{q}\mathbf{q}'\mathbf{q}''}N\,\delta_{\mathbf{q}+\mathbf{q}'+\mathbf{q}'',\,\mathbf{g}}\,\mathsf{A}(\mathbf{q}',\mathbf{q}'')\colon \mathbf{U}_{\mathbf{q}}\mathbf{U}_{\mathbf{q}'}\mathbf{U}_{\mathbf{q}''}.
\tag{2.127}
\end{align*}
我们看到，模间耦合系数的张量就是非谐项张量的一个 Fourier 变换。尤其值得注意的是，对 $\mathbf{l}$ 的求和引入了一个形如 (2.92) 的因子，只有当
\begin{equation*}
\mathbf{q}+\mathbf{q}'+\mathbf{q}''=\mathbf{g},
\tag{2.128}
\end{equation*}
时它才不为零，这里 $\mathbf{g}$ 照例是倒格矢或零。于是，相互耦合的各个过程，其晶体动量就要满足一个条件。

我们可以这样来理解。假设我们把非谐项当作对晶格模运动的微扰，这些运动受方程 (2.11) 支配。通过一个正则变换，这些方程可以约化为一组独立简谐振子的运动方程。众所周知，这类方程在量子力学中用一个哈密顿量来表示，而这个哈密顿量又可以再变换成一个含有各模中量子的产生算符与湮没算符的表达式。

事实上，我们发现可以写出
\begin{equation*}
\mathbf{U}_{\mathbf{q}}=-i\left(\frac{\hbar}{2NM\nu_{\mathbf{q}}}\right)^{1/2}\mathbf{e}_{\mathbf{q}}\,(a_{\mathbf{q}}-a^{*}_{-\mathbf{q}}),
\tag{2.129}
\end{equation*}
其中 $\mathbf{e}_{\mathbf{q}}$ 是谱中某一支内波矢为 $\mathbf{q}$ 的模在偏振方向上的单位矢量，而 $a_{\mathbf{q}}$、$a^{*}_{\mathbf{q}}$ 是该模中一个声子的湮没算符与产生算符。当然，偏振矢量由本征值方程 (2.13) 对频率 $\nu_{\mathbf{q}}$ 的解来确定。

用这些算符来表示，由谐和项导出的哈密顿量约化成非常简单的形式
\begin{equation*}
\mathcal{H}_0=\sum_{\text{偏振}}\sum_{\mathbf{q}}\hbar\nu_{\mathbf{q}}\,(a_{\mathbf{q}}a^{*}_{\mathbf{q}}+a^{*}_{\mathbf{q}}a_{\mathbf{q}}),
\tag{2.130}
\end{equation*}
其本征态就是声子态，譬如 $|n_{\mathbf{q}}\rangle$，即第 $\mathbf{q}$ 个模中有 $n_{\mathbf{q}}$ 个量子。把 (2.129) 代入非谐项 (2.127)，就得到一个非常复杂的表达式，其中含有算符 $a_{\mathbf{q}}$、$a_{\mathbf{q}'}$ 等的乘积。例如，有一项包含（在挪动了一些正负号之后）
\begin{equation*}
\delta_{\mathbf{q}-\mathbf{q}'-\mathbf{q}'',\,\mathbf{g}}\,a^{*}_{\mathbf{q}}a_{\mathbf{q}'}a_{\mathbf{q}''}.
\tag{2.131}
\end{equation*}
当把这一项作为微扰作用到 $\mathcal{H}_0$ 的本征态上时，会引起这样的跃迁：在态 $\mathbf{q}$ 中产生一个声子，同时在态 $\mathbf{q}'$ 和 $\mathbf{q}''$ 中各消灭一个声子。相应的选择定则是
\begin{equation*}
\mathbf{q}'+\mathbf{q}''=\mathbf{q}-\mathbf{g}.
\tag{2.132}
\end{equation*}

这与式 (2.100) 相类似。若 $\mathbf{g}$ 为零，我们就说晶体动量守恒：动量 $\hbar\mathbf{q}'$ 和 $\hbar\mathbf{q}''$ 被消灭，凑出了动量 $\hbar\mathbf{q}$。当 $\mathbf{g}$ 不为零时，晶体动量不再守恒——但它只能改变 $\hbar\mathbf{g}$ 这么多。于是，\emph{声子–声子相互作用}可以像外束 X 射线的散射那样，用 N 过程和倒逆过程（U 过程）来描述。显然还有一些别的过程，其中一个声子被消灭、两个声子被产生——即 (2.131) 的共轭过程（译注：原文印作 (2.130)）。

可以证明，真实过程必须满足能量守恒；对 (2.132) 的情形，必须有
\begin{equation*}
\hbar\nu_{\mathbf{q}'}+\hbar\nu_{\mathbf{q}''}=\hbar\nu_{\mathbf{q}}.
\tag{2.133}
\end{equation*}
这一论证的依据是晶格模的时间依赖性。

显然，声子–声子相互作用实际的系数非常复杂：它们由非谐系数导出，还同偏振矢量等混杂在一起。不同偏振的模之间真实跃迁的规则也很复杂，因为这类跃迁要求能量条件与晶体动量条件同时得到满足。同样清楚的是，还可以有更高阶的声子–声子相互作用，例如 $n$ 个声子被消灭而 $n'$ 个声子被产生。这类项既可能来自晶格势能 Taylor 展开中的 $n+n'$ 阶项，也可能来自微扰级数中的更高阶项。它们必须满足类似 (2.128) 的选择定则——晶体动量的总改变必须是 $\hbar\mathbf{g}$，这里 $\mathbf{g}$ 为零或倒格矢。声子–声子相互作用过程在晶格热传导理论中十分重要，我们将在 §7.10 再回到这一点。

对于\emph{量子晶体}（quantum crystals）——特别是固态氦——上述晶格动力学的研究路线需要修正；在这类晶体中，原子的零点运动（zero-point motion）与晶格间距相当。由一串阶数越来越高的系数所定义、相继描述越来越复杂的声子相互作用过程的含时微扰级数，不再绝对收敛：我们甚至会发现，在平衡晶格间距处对势能求导而算得的二阶谐和力常数 (2.4) 取\emph{负}值。要修补这套代数，我们必须求助于\emph{多体理论}（many-body theory）的标准技巧之一。例如，我们可以走一走\emph{图形方法或簇展开}（diagrams or cluster expansions）的路子，在把微扰级数中各种无穷多项的组合求和之后，对所有相互作用系数进行重整化。或者，循着更直接的物理直觉，仿照 Hartree 方法的精神，为原子设立试探波函数，并对参数寻找变分条件。两种做法的结果大体相同。系统的行为相当正常：有声子谱、有声子–声子相互作用等等，只是有效力常数和更高阶的系数必须\emph{自洽地}（self-consistently）确定。于是，应把 (2.4) 换成大致如下的形式
\begin{equation*}
\mathsf{G}_{ll'}=\left\langle\left|\,\frac{\partial^2\mathcal{V}}{\partial\mathbf{u}_l\,\partial\mathbf{u}_{l'}}\,\right|\right\rangle,
\tag{2.134}
\end{equation*}
其中势能对原子位移的导数，要在某个与力常数的这些取值相一致的声子态 $|\;\rangle$ 中作为\emph{期望值}（expectation value）来求值。

\section{非完整格子的振动}

真实的晶体绝不可能是完美的，也不可能是无限延伸的；我们必须考虑缺陷和表面对晶格振动的影响。这个课题是没有尽头的，会通向诸如完全无序系统（如液体和玻璃）的动力学这类棘手的问题。不过，如果我们只限于讨论孤立的缺陷——例如质量不同的替位原子（substitutional atoms）——那么如今已能理解其中不少有趣的物理现象。

论证的要点通常被淹没在一大堆琐碎的细节之中。为了降低代数运算表面上的复杂性，让我们采用一种抽象的矩阵记法；例如，\emph{完美的}（Bravais）格子振动的一般方程在这种记法下便写成
\begin{equation*}
\mathsf{L}(\nu^2)\,u=0
\tag{2.135}
\end{equation*}
它代表诸如
\begin{equation*}
M\nu^2\mathbf{u}_l-\sum_{l'}\mathsf{G}_{ll'}\cdot\mathbf{u}_{l'}=0.
\tag{2.136}
\end{equation*}
这些方程本质上与 (2.5) 相同；列矩阵 $u$ 代表频率为 $\nu$ 的模中各位移 $\mathbf{u}_l$ 的全体，此时原子间相互作用力为 $\mathsf{G}_{ll'}$。动力学矩阵（dynamical matrix）$\mathsf{L}$ 具有晶格的平移对称性。按 §2.1 的论证，在每个声子频率 $\nu=\nu_{\mathbf{q}}$ 处 (2.135) 都有一个根：因为这时我们能为相应的声子态求得一个（归一化的）偏振矢量 $\mathbf{U}_{\mathbf{q}}$
\begin{equation*}
\mathbf{u}_l=\mathbf{U}_{\mathbf{q}}\,e^{i\mathbf{q}\cdot\mathbf{l}}
\tag{2.137}
\end{equation*}
而这只需解本征矢方程 (2.13) 即可。

现在假设我们改变了 $\mathbf{l}=0$ 处（其实是任选的一个格点）原子的质量，改变量为 $\delta M$，并且（或者）改变了该原子与晶体中近邻的相互作用，改变量为 $\delta\mathsf{G}_{ll'}$。我们的动力学方程 (2.136) 这时就应读作
\begin{equation*}
M\nu^2\mathbf{u}_l-\sum_{l'}\mathsf{G}_{ll'}\cdot\mathbf{u}_{l'}+\delta M\,\nu^2\mathbf{u}_0+\sum_{l'}\delta\mathsf{G}_{ll'}\cdot\mathbf{u}_{l'}=0,
\tag{2.138}
\end{equation*}
我们把它符号式地记为
\begin{equation*}
\mathsf{L}u+\delta\mathsf{L}u=0.
\tag{2.139}
\end{equation*}
这样，物理系统的这些改变就可以表示为原先完美晶体的动力学矩阵的一种微扰。

在线性链的特殊情形（参照 (2.15)），这些方程多少可以直接看出解来。但 Green 函数法（Green function method）适用面更广，而且也复杂不了多少。诀窍是把 (2.138) 改写成代数上等价的形式
\begin{equation*}
(\mathbf{1}+\mathsf{R}\,\delta\mathsf{L})\,u=0,
\tag{2.140}
\end{equation*}
其中预解式（resolvent）即 Green 函数 $\mathsf{R}$ 是动力学矩阵 $\mathsf{L}$ 的逆矩阵。只要 $\mathsf{R}$ 存在而且算得出来，我们就有了解决 (2.140) 所代表的线性方程的头绪。

预解式必须是 $\nu^2$ 的函数，满足抽象关系式
\begin{equation*}
\mathsf{R}(\nu^2)\,\mathsf{L}(\nu^2)=\mathbf{1}.
\tag{2.141}
\end{equation*}
很容易验证，本征矢 (2.137) 恰好给出这一问题的一组完全解：利用 (2.13) 诸解的标准正交性等性质，我们发现矩阵 $\mathsf{R}(\nu^2)$ 可以表示为
\begin{equation*}
\mathsf{R}_{ll'}(\nu^2)=\frac{1}{NM}\sum_{\mathbf{q}}\frac{\mathbf{U}^{*}_{\mathbf{q}}\mathbf{U}_{\mathbf{q}}}{\nu^2-\nu^2_{\mathbf{q}}}\,e^{i\mathbf{q}\cdot(\mathbf{l}-\mathbf{l}')}.
\tag{2.142}
\end{equation*}
这里我们同样出于简单，假定对所有波矢 $\mathbf{q}$ 的求和已扩展到把声子偏振的各种模式都包括进来。这个公式还使用了并矢（dyadic）约定，以顾及动力学矩阵的笛卡尔张量特征；不过，这仍然只是我们本应为例如一维链写下的那类公式的一个相当平凡的推广。

原则上，我们可以对 $(\mathbf{l}-\mathbf{l}')$ 的一切取值算出 (2.142)。实际上这并不必要，因为我们假定微扰 $\delta\mathsf{L}$ 高度局域（localized）在缺陷格点的邻域内。例如，考虑同位素杂质（isotopic impurity）的情形，此时 (2.138) 中只有质量发生了变化。把 (2.142) 代入 (2.140)，便得到一个只含 $\mathbf{u}_0$——即该格点上的矢量位移——的方程。其中只出现 $\mathsf{R}_{00}(\nu^2)$ 的求和，即
\begin{equation*}
\Bigl\{\mathbf{1}+\frac{\delta M}{M}\,\frac{\nu^2}{N}\sum_{\mathbf{q}}\frac{\mathbf{U}^{*}_{\mathbf{q}}\mathbf{U}_{\mathbf{q}}}{\nu^2-\nu^2_{\mathbf{q}}}\Bigr\}\cdot\mathbf{u}_0=0.
\tag{2.143}
\end{equation*}
笛卡尔张量 $\{\;\}$ 的 $3\times3$ 行列式为零，这一条件加在 $\nu^2$ 上，从而给出非完整晶格某一简正模的振动频率。更复杂的情形——几个格点上的力常数都有变化——显然会给出类似的方程，只是行列式的阶数更高。

遗憾的是，这些方程很难求解：即便是同位素替代这种初等情形，算起来也十分繁复。但为了理解究竟会发生什么，让我们撇开诸如 $\mathbf{u}_0$ 的矢量性这类几何特征，把 (2.143) 当作一个标量方程来处理。简正模的频率是如下形式方程的根
\begin{equation*}
1+\left(\frac{\delta M}{M}\right)\frac{1}{N}\sum_{\mathbf{q}}\frac{\nu^2}{\nu^2-\nu^2_{\mathbf{q}}}=0.
\tag{2.144}
\end{equation*}

\begin{figure}[H]
\centering
\includegraphics[width=0.72\textwidth]{fig_34.pdf}
\caption*{图 34\quad 对含同位素杂质的晶格的模，式 (2.143) 的图解求法。}
\end{figure}

这些根可以用作图法求得（图 34）。我们要找的是函数
\begin{equation*}
f(\nu^2)\equiv\frac{1}{N}\sum_{\mathbf{q}}\frac{\nu^2}{\nu^2-\nu^2_{\mathbf{q}}}
\tag{2.145}
\end{equation*}
与水平线 $(-M/\delta M)$ 相交的那些点。

例如，设 $\delta M$ 为负。(2.144) 的每个根必定位于 $f(\nu^2)$ 的某个极点 $\nu^2_{\mathbf{q}}$ 的上方；受扰系统的各简正模在频率上与完美晶体的各简正模交错排布。由于 $\nu_{\mathbf{q}}$ 的取值构成一个稠密的\emph{带}，其间隔像 $1/N$ 那样趋于零，新解中的大多数与原有的解无从区分。但是再看一眼图 34 就会注意到，最高的那个根并不受这种约束，它可以离开带的顶端 $\nu_{\max}$ 一段有限的距离。当 $\nu$ 大于 $\nu_{\max}$ 时，求和 (2.145) 可以近似地用对谱密度（spectral density）(2.65) 的一个积分来表示，即
\begin{equation*}
\bar{f}(\nu^2)=\int^{\nu^{\max}}\frac{\nu^2\mathcal{D}(\nu')}{\nu^2-\nu'^2}\,d\nu'.
\tag{2.146}
\end{equation*}
在这个函数等于 $(-M/\delta M)$ 的频率处，就出现一个与缺陷相联系的特殊模。

这个特殊模必定是\emph{局域}的。这来自一条普遍定理：\emph{晶格的振荡运动在简正模频谱之外的频率处，其波数为虚数}。以 §2.2 的线性链模型为例。若把 $q=i\kappa$ 代入 (2.16)，我们便解出了 (2.15)，其频率由 (2.18) 的类比式给出，即
\begin{equation*}
\nu_{\kappa}=\sqrt{\left(\frac{\alpha}{M}\right)}\;2\sinh\left(\frac{\kappa a}{2}\right),
\tag{2.147}
\end{equation*}
对 $\nu_{\kappa}$ 的任何取值它都能得到满足。当然，作为完美晶格的简正模，这个解析解在物理上是不允许的，因为位移振幅 $u_l$ 会沿链像 $\exp(\pm\kappa l)$ 那样指数式地增长或衰减，从而不能满足循环边界条件。然而在非完整晶体中，可以把缺陷邻域的条件安排成使振动振幅向四面八方衰减，于是在样品远处的边界上变得可以忽略。从物理上讲，图 34 中的杂质原子比晶体中周围的原子来得轻，它作独立“Einstein”振荡时的固有频率会比较高。观察到的就是位于扩展声子模带之上的一个简正模，其中只有少数几个近邻原子随杂质协同运动（图 35(a)）。

缺陷位置上局域振动的其他种种例子都可以类似地理解。例如，如果我们把某个原子周围的弹簧常量加硬一些，就会观察到同样的效应。另一方面，比平常更重一些的原子具有较低的 Einstein 频率，因此我们应当能看到一个局域模从某个正常频带之\emph{下}脱离出来——例如从双原子晶体的光学支中脱离出来。

说来也怪，表面模（surface modes）的理论竟循着同样的路数。当然，我们可以从经典弹性理论中搬用关于连续介质自由表面上 Rayleigh 波与 Love 波的标准讨论，或许还要对表面平面内波矢的大小加上一个 Debye 限制。但这些波对应于那些局域在表面上、并朝晶体内部指数式衰减的振动模。虽然离散晶格中相应的问题看上去非常复杂，

\begin{figure}[H]
\centering
\includegraphics[width=0.72\textwidth]{fig_35.pdf}
\caption*{图 35\quad (a) 轻杂质在声子带之上的局域模。(b) 重杂质在声子带之内的共振模。}
\end{figure}

但若把系统处理成一列晶格平面——它们从表面平面处的一个重大缺陷（例如劲度为零的弹簧常量）出发，沿一维向晶体内部延伸——还是能学到许多东西。这样一来，这类模应当出现在双原子晶体中、声子谱的光学支与声学支之间的带隙里。表面模可能对微小晶粒的比热有可检测的贡献，并且在某些电子器件中很重要。

局域杂质模的特征频率可以通过红外波段的光吸收直接观测（见 §8.4）。而且，有意思的是，如果杂质的固有振动频率落在某个声子带之\emph{内}，也会产生一条有所展宽的吸收线——例如在 (2.143) 的简单同位素情形中取 $\delta M$ 为正。这是\emph{共振}（resonance）的一个例子；在峰频率处，振动振幅在杂质格点上要比晶体其余部分大得多，但它并不随距离指数式地消衰，因此不是严格局域的（图 35(b)）。不出所料，在这个频率上杂质对入射声子的散射也非常强烈。

事实上，共振的理论只不过是局域态理论的一种解析延拓。共振频率是 (2.144) 的解，只是把其中的求和 $f(\nu^2)$ 换成 $\bar{f}(\nu^2)$ 的积分 (2.146) 的主部；而共振的\emph{宽度}则由同一积分的虚部给出，只不过现在要在 $\nu<\nu_{\max}$ 的情形下取值。这类关系的证明通常是在“球对称势中量子力学粒子的 Schrödinger 方程之解”的语境下给出的，但这一论证普遍适用于连续介质和规则晶格中的波。
